% concatenated from iclr_25.tex, src/math-commands.tex, src/intro_danilo_edit.tex, src/background.tex, src/theory.tex, src/non-linear.tex, src/experiments.tex, src/conclusion.tex, src/appendix.tex

\documentclass{article} % For LaTeX2e
\usepackage{iclr2025_conference,times}

\usepackage{amsmath}
\usepackage{amsfonts}
\usepackage{aligned-overset}
\usepackage{amsthm}
\usepackage{amssymb}
\usepackage{nicefrac}
\usepackage[ruled,vlined]{algorithm2e}
\usepackage{caption}
\usepackage{subcaption}
\usepackage{wrapfig}
\usepackage{enumitem}

\newtheorem{theorem}{Theorem}[section]
\newtheorem{corollary}[theorem]{Corollary}
\newtheorem{lemma}[theorem]{Lemma}
\newtheorem{proposition}[theorem]{Proposition}
\newtheorem{algo}[theorem]{Algorithm}
\theoremstyle{definition}
\newtheorem{definition}[theorem]{Definition}
\newtheorem{assumption}[theorem]{Assumption}
\newtheorem{remark}[theorem]{Remark}


% Optional math commands from https://github.com/goodfeli/dlbook_notation.
% using our own file
\newcommand{\reals}{\mathbb{R}}
\newcommand*\closure[1]{\operatorname{cl}#1}
\newcommand{\nnumbers}{\mathbb{N}}
\newcommand{\vi}{\operatorname{VI}}
\newcommand{\ri}{\operatorname{ri}}
\newcommand{\diag}{\operatorname{diag}}
\newcommand{\ran}{\operatorname{ran}}
\newcommand{\sigmoid}{\operatorname{sigmoid}}
\newcommand{\softmax}{\operatorname{softmax}}
\newcommand{\celu}{\operatorname{CELU}}
\newcommand{\inner}[2]{\langle #1, #2 \rangle}
\newcommand{\prox}{\operatorname{Prox}}
\newcommand{\id}{\operatorname{Id}}
\newcommand{\sol}{\operatorname{sol}}
\newcommand{\zer}{\operatorname{zer}}
\newcommand{\fixPoint}{\operatorname{fix}}
\newcommand{\eig}{\operatorname{eig}}
\newcommand{\expect}[2][]{\mathbb{E}_{#1}\left[ #2\right]}
% \newtheorem{algo}[theorem]{Algorithm}

% \newtheorem{example}[theorem]{Example}


\usepackage{hyperref}
\usepackage{url}




\title{Solving hidden monotone variational inequalities with surrogate losses}
% Authors must not appear in the submitted version. They should be hidden
% as long as the \iclrfinalcopy macro remains commented out below.
% Non-anonymous submissions will be rejected without review.
\author{Ryan D'Orazio, Danilo Vucetic \& Zichu Liu \\
Mila Qu\'ebec AI Institute, Universit\'e de Montr\'eal\\
Montr\'eal, QC, Canada \\
\texttt{\{ryan.dorazio, danilo.vucetic, zichu.liu\}@mila.quebec} \\
\And
Junhyung Lyle Kim  \\
Department of Computer Science, Rice University \\
Houston, TX, USA \\
\texttt{jlylekim@rice.edu}\\
\AND
Ioannis Mitliagkas \& Gauthier Gidel \\
Mila Qu\'ebec AI Institute, Universit\'e de Montr\'eal, CIFAR AI Chair\\
Montr\'eal, QC, Canada \\
\texttt{\{ioannis, gidelgau\}@mila.quebec}
}
% \author{Antiquus S.~Hippocampus, Natalia Cerebro \& Amelie P. Amygdale \thanks{ Use footnote for providing further information
% about author (webpage, alternative address)---\emph{not} for acknowledging
% funding agencies.  Funding acknowledgements go at the end of the paper.} \\
% Department of Computer Science\\
% Cranberry-Lemon University\\
% Pittsburgh, PA 15213, USA \\
% \texttt{\{hippo,brain,jen\}@cs.cranberry-lemon.edu} \\
% \And
% Ji Q. Ren \& Yevgeny LeNet \\
% Department of Computational Neuroscience \\
% University of the Witwatersrand \\
% Joburg, South Africa \\
% \texttt{\{robot,net\}@wits.ac.za} \\
% \AND
% Coauthor \\
% Affiliation \\
% Address \\
% \texttt{email}
% }

% The \author macro works with any number of authors. There are two commands
% used to separate the names and addresses of multiple authors: \And and \AND.
%
% Using \And between authors leaves it to \LaTeX{} to determine where to break
% the lines. Using \AND forces a linebreak at that point. So, if \LaTeX{}
% puts 3 of 4 authors names on the first line, and the last on the second
% line, try using \AND instead of \And before the third author name.

\newcommand{\fix}{\marginpar{FIX}}
\newcommand{\new}{\marginpar{NEW}}

\iclrfinalcopy % Uncomment for camera-ready version, but NOT for submission.
\begin{document}


\maketitle


\begin{abstract}
Deep learning has proven to be effective in a wide variety of loss minimization problems.
However, many applications of interest, like minimizing projected Bellman error and min-max optimization, cannot be modelled as minimizing a scalar loss function but instead correspond to solving a variational inequality (VI) problem.
This difference in setting has caused many practical challenges as naive gradient-based approaches from supervised learning tend to diverge and cycle in the VI case.
In this work, we propose a principled surrogate-based approach compatible with deep learning to solve VIs.
We show that our surrogate-based approach has three main benefits: (1) under assumptions that are realistic in practice (when hidden monotone structure is present, interpolation, and sufficient optimization of the surrogates), it guarantees convergence, (2) it provides a unifying perspective of existing methods, and (3) is amenable to existing deep learning optimizers like ADAM.
Experimentally, we demonstrate our surrogate-based approach is effective in min-max optimization and minimizing projected Bellman error. Furthermore, in the deep reinforcement learning case, we propose a novel variant of TD(0) which is more compute and sample efficient.
\end{abstract}

\section{Introduction}

Most machine learning approaches learn from data by minimizing a loss function with respect to model parameters.
Despite the lack of global convergence guarantees in deep learning, losses can often still be minimized with an appropriately tuned first-order adaptive method such as ADAM~\citep{kingma2014adam}.
Unfortunately, outside of scalar loss minimization, deep learning becomes more challenging: the dynamics of variational inequality (VI) problems (\textit{e.g.}, min-max) are often plagued with rotations and posses no efficient stationary point guarantees~\citep{daskalakis2021complexity}.
Thus, the additional challenges posed by VI problems do not allow one to easily plug in existing techniques from deep learning.

We propose a practical and provably convergent algorithm for solving VI problems in deep learning.
Importantly, our approach is compatible with any black-box optimizer, including ADAM and its variants~\citep{kingma2014adam}, so long as they are effective at descending a scalar loss.
Our approach guarantees convergence by leveraging the hidden structure commonly found in machine learning.
More precisely, many applications admit a hidden structure corresponding to the following problem: 
\begin{align}\label{eq:vi}
\text{find }z_\ast \text{ such that }
    \inner{F(z_\ast)}{z-z_\ast} \geq 0 \quad \forall z \in \mathcal{Z} = \closure{\{g(\theta) : \theta \in \reals^d\}},
\end{align}
where the mapping $g: \reals^d \to \mathcal{Z} \subseteq \reals^n$ maps model parameters to model outputs, and the set $\mathcal{Z}$ encodes the closure of the set of realizable outputs from the chosen model.
Since $F$ is defined over model outputs, it will often be structured, e.g., monotone and Lipschitz.

% We propose a practical and provably convergent algorithm for solving VI problems in deep learning.
% Our approach reduces VI problems to the approximate minimization of a sequence of surrogate losses, we effectively translate VI problems into scalar minimization problems, which are much easier to solve. We ensure the convergence of this method with respect to a novel $\alpha$-descent condition by leveraging the monotonicity and smoothness of common loss functions with respect to model outputs.
% % By leveraging the monotonicity and smoothness of common loss functions with respect to model outputs to ensure convergence with respect to a novel $\alpha$-descent condition.
% Importantly, our method allows for any black-box optimizer to be used, including ADAM and its variants~\citep{kingma2014adam}.
% More precisely, many loss functions admit a hidden structure corresponding to the following VI problem: find $z_\ast$ such that
% \begin{align}\label{eq:vi}
%     \inner{F(z_\ast)}{z-z_\ast} \geq 0 \quad \forall z \in \mathcal{Z} = \closure{\{g(\theta) : \theta \in \reals^d\}},
% \end{align}
% where the mapping $g: \reals^d \to \mathcal{Z} \subseteq \reals^n$ maps model parameters to model outputs, and the set $\mathcal{Z}$ encodes the closure of the set of realizable outputs from the chosen model.
% Since $F$ is defined over model outputs, it will often be monotone and smooth.
% %\footnote{In the scalar minimization case, $F$ is the gradient of the loss function: $\nabla \ell$.}
% % In other words: a convex problem is often hiding in plain sight!

In order to leverage the hidden structure in VI problems, we employ surrogate losses, which have been studied in the scalar minimization case by \citet{pmlr-v119-johnson20b} and \citet{vaswani2021general}, among others.
The surrogate loss approach has been shown to be scalable with deep learning and has been used in modern reinforcement learning policy gradient methods \citep{schulman2015trust,schulman2017proximal,abdolmaleki2018maximum}.
Moreover, these surrogate loss approaches have been shown to improve data efficiency in cases where evaluating $F$ is expensive like interacting with a simulator in reinforcement learning~\citep{vaswani2021general,lavington2023target}.
% Crucially, surrogate losses have long been used in policy gradient methods such as: TRPO~\citep{schulman2015trust}, PPO~\citep{schulman2017proximal}, and MPO~\citep{abdolmaleki2018maximum}.
% Although these policy gradient methods weren't originally designed to take advantage of hidden structure, they can be naturally interpreted as performing an approximate gradient step in the space of model outputs~\citep{pmlr-v119-johnson20b, vaswani2021general}.
% More importantly, these surrogate loss approaches are scalable with deep learning and allow for efficient use of data in cases where evaluating $F$ is expensive like interacting with a simulator in reinforcement learning or imitation learning~\citep{vaswani2021general,lavington2023target}.

Our approach reduces the VI in \eqref{eq:vi} to the approximate minimization of a sequence of surrogate losses $\{\ell_t\}_{t \in \nnumbers}$, which are then used to produce a sequence of parameters $\{\theta_t\}_{t\in \nnumbers}$.
To ensure convergence we propose a new $\alpha$-descent condition on $\ell_t$, allowing for a dynamic inner-loop that makes no assumption on how the surrogate losses are minimized, thereby allowing for any deep-learning optimizer to minimize the scalar loss $\ell_t$.
With our $\alpha$-descent condition we provide convergence guarantees to a solution in the space of model predictions $\{z_t = g(\theta_t)\}_{t \in \nnumbers} \to z_\ast$ for a sufficiently small $\alpha$.
Our general method is summarized in Algorithm \ref{alg:surr}.
% where the practicality of our method is evident: one simply needs to define their VI operator $F$ and their loss function $\ell_t$ to solve a VI problem.
\vspace{-3mm}
\begin{algorithm}
\caption{$\alpha$-descent on surrogate}\label{alg:surr}
 \SetAlgoLined
 \SetAlgoNoEnd
  \KwIn{Outer loop interactions $T$, initial parameters $\theta_1 \in \reals^d$, step size $\eta$ for surrogate loss, $\alpha \in (0,1)$, optimizer update $\mathcal{A}: \mathcal{L} \times \reals^d \to \reals^d$.}
  \For{$t = 1 \leftarrow$ \KwTo $T$}{
    Compute VI operator: $F(g(\theta_t))$\\
    Set surrogate loss: $\ell_t(\theta) = \frac{1}{2}\|g(\theta) - (g(\theta_t)-\eta F(g(\theta_t))\|^2$\\
    $\theta_s \leftarrow \theta_t$\\
    \While{$\ell_t(\theta_s) -\ell_t^\ast > \alpha^2(\ell_t(\theta_t) - \ell_t^\ast)$}{
    Update parameters with optimizer:
    $\theta_s \leftarrow \mathcal{A}(\ell_t, \theta_s)$
    }
    $\theta_{t+1} \leftarrow \theta_s$
  } 
  \Return{$\theta_{T+1}$}
\end{algorithm}
\vspace{-6mm}
% \textcolor{blue}{JLK: we should explain the algorithm a bit here}


In summary, we propose a novel algorithm to solve VI problems by exploiting the hidden structure of common loss functions. We demonstrate the efficacy of our method with a variety of experiments, which themselves make further connections and innovations on prior work. Our specific contributions are as follow:
\vspace{-8pt}
\begin{itemize}[leftmargin=*]
\setlength\itemsep{0pt}
    \item We provide the \textbf{first extension of surrogate losses to VI problems}.
    We also show a clear separation in difficulty of using surrogate methods in VI problems when compared to scalar minimization; specifically, we show that divergence is possible in a strongly monotone VI problem where in contrast convergence is guaranteed in the non-convex scalar minimization case.
    % , divergence is possible in a strongly monotone VI problem. 
        % We also show that VI problems are significantly more difficult to solve than scalar minimization problems, and may diverge even if consistent progress is made on the sequence surrogate losses.
    \item We propose the \textbf{$\mathbf{\alpha}$-descent condition with convergence guarantees}.
    This condition allows for global convergence while avoiding common assumptions such as enforcing errors to be summable or globally upper-bounded.
        % The $\alpha$-descent condition is used to quantify progress on the sequence of surrogate losses, allowing for global convergence guarantees without sweeping assumptions like bounded losses.
        % This condition allows for global convergence without requiring sweeping assumptions like bounded losses.
        % We also demonstrate a connection with biased gradients in VI problems, thereby providing new insights and results on the difficulties of solving VIs with bias.
        %forcing a global upper bound on all the surrogate losses at each iteration and better models implementations found in practice where a fixed number of gradient descent steps are used to minimize $\{\ell_t\}_{t \in \nnumbers}$.
        % Our analysis is mainly provided in the deterministic setting with preliminary results in the stochastic case.
        % In contrast to the scalar minimization case where any $\alpha \in [0,1)$ is sufficient for convergence, 
        % we demonstrate the inherent difficulties of extending surrogate losses to VI problems by showing divergence is possible with $\alpha <1 $.
        % Moreover, our theoretical results can be viewed within the context of solving VIs with biased gradients. 
        % We therefore provide new insights and results on the difficulties of solving VIs with bias.

    \item \textbf{Unifying perspective of pre-conditioning methods}.
        With our surrogate loss approach, we unify existing pre-conditioning methods ~\citep{bertsekas2009projected,mladenovic2022generalized,sakos2024exploiting} by showing they are equivalent to using the Gauss-Newton method as the optimizer $\mathcal{A}$ in Algorithm \ref{alg:surr} to minimize the surrogate losses.
        We demonstrate the value of this new perspective by providing natural extensions of their methods with better empirical robustness.

    \item \textbf{Experimental results and new TD variants}. We demonstrate the performance and versatility of surrogate loss-based optimization in a variety of VI problems. In Sections \ref{sec:experiments_minmax} and \ref{sec:experiments_PBE} we complete experiments in min-max optimization and value prediction tasks, respectively. Importantly, we propose a new data-efficient variant to TD(0) which significantly outperforms prior approaches.
    %In Section \ref{sec:experiments_PBE}, we perform value prediction tasks in reinforcement learning and propose a new, more data-efficient, variant to temporal difference (TD) learning.
\end{itemize}


%is comprised of a number of physics-based experiments whereby the value of an expert policy was learned \textcolor{red}{via successive TD updates.} Our value prediction experiments demonstrate that the surrogate approach outlined here is capable of learning complex, non-stationary value functions, while ... 

% We study a surrogate loss framework that leverages the hidden convexity often present in modern machine learning applications including min-max optimization and problems that cannot be modeled by minimizing a scalar loss function. 
% In most applications, such convexity (or monotonicity, more generally) is present with respect to model predictions but not with respect to the model's parameters.

% Therefore the problem in prediction space can be often viewed as a well-conditioned constrained variational inequality (VI) problem, where the constraint encodes the model capacity.
% More precisely, we are interested in solving the following variational inequality problem
% \begin{align}\label{eq:vi}
%     \inner{F(z_\ast)}{z-z_\ast} \geq 0 \quad \forall z \in \mathcal{Z} = \closure{\{g(\theta) : \theta \in \reals^d\}}.
% \end{align}
% Where the mapping $g: \reals^d \to \mathcal{Z} \subseteq \reals^n$ encodes the the set of realizable outputs $\mathcal{Z}$ from function approximation like deep learning models.

% Our framework leverages the structure in prediction space by selecting a sequences of losses $\{\ell_t\}_{t \in \nnumbers}$ for which are then used to produce a sequence of parameters $\{\theta_t\}_{t\in \nnumbers}$.
% We propose an $\alpha$ descent condition on $\ell_t$ that guarantees convergence to a solution in the space of model predictions $\{z_t = g(\theta_t)\}_{t \in \nnumbers} \to z_\ast$ for a sufficiently small $\alpha$.
% More importantly, our approach is flexible and compatible with deep learning allowing for the use of any black-box optimizers including gradient descent.



% \paragraph{Extension of surrogate losses to VI problems and challenges.}
% This work provides the first extension of surrogate losses to VI problems.
% Although our approach includes scalar minimization as a special case, we show there is a clear separation in problem difficulty.
% In the scalar minimization case, if $F$ is smooth then any progress on the surrogate loss is sufficient for convergence~\citep{vaswani2021general,lavington2023target}; meanwhile, we show that it is possible to diverge in the VI case with consistent progress on the surrogate losses even if $F$ is strongly monotone.

% \paragraph{Inner-loop condition and convergence guarantees.}
% In contrast to existing approaches, we control our inner-loop via an $\alpha$-descent condition.
% This condition allows for global convergence without forcing a global upper bound on all the surrogate losses at each iteration and better models implementations found in practice where a fixed number of gradient descent steps are used to minimize $\{\ell_t\}_{t \in \nnumbers}$.
% Our analysis is mainly provided in the deterministic  setting with preliminary results in the stochastic case.
% % In contrast to the scalar minimization case where any $\alpha \in [0,1)$ is sufficient for convergence, 
% % we demonstrate the inherent difficulties of extending surrogate losses to VI problems by showing divergence is possible with $\alpha <1 $.
% Moreover, our theoretical results can be viewed within the context of solving VIs with biased gradients. 
% We therefore provide new insights and results on the difficulties of solving VIs with bias.

% \paragraph{Unifying perspective of pre-conditioning methods.}
% With our surrogate loss approach we are able to unify existing pre-conditioning methods ~\citep{bertsekas2009projected,mladenovic2022generalized,sakos2024exploiting} by showing they are equivalent to picking the Gauss-Newton method as the optimizer $\mathcal{A}$ in Algorithm \ref{alg:surr} to minimize the surrogate losses.
% Additionally, we demonstrate the value of this new prospective by providing natural extensions of their methods with better empirical robustness and performance.

% Additionally, we demonstrate the value of this new prospective by providing a guarantee with respect to our $\alpha$-descent condition for these methods.
% \paragraph{Experiments.}
% We complement our analysis with  experiments in min-max optimization and minimizing projected Bellman error (PBE) in reinforcement learning (RL).
% We also provide a new variant of temporal difference learning (TD), a generalization of TD(0) that is more data efficient.





\section{Background and Related Work}
\paragraph{Notation.} We use $\inner{x}{y} = \sum_{i=1}^n x^i y^i$ to denote the standard inner product over $\reals^n$ and $\|x\| = \sqrt{\inner{x}{x}}$ to be the Euclidean norm. 
We write $\|x\|^2_{\Xi}$ to mean $\inner{x}{\Xi x}$,  and $\|x\|_{\Xi}$ is a norm if and only if $\Xi$ is positive definite. 
For a set $\mathcal{X}$ we denote $\closure{\mathcal{X}}$ its closure and $\ri \mathcal{X}$ its relative interior.
For a given set, which will be clear from context, we denote $\Pi(x)$ as the Euclidean projection of $x$ onto the set and similarly $\Pi_{\Xi}(x)$ the projection with respect to $\|x\|_{\Xi}$. 
We use $\id$ to denote the identity matrix.
A matrix $A$ has lower and upper bounded singular values if
there exists $\sigma_{\min},\sigma_{\max} \in (0,\infty)$ such that for any $x$ we have $\sigma_{\min}^2\|x\|^2 \leq \inner{x}{A^\top A x} \leq \sigma_{\max}^2\|x\|^2$.
If a matrix $A$ is invertible we write $A^{-1}$ otherwise we denote the pseudo-inverse as $A^{\dag}$.
For a given function $g: \reals^d \to \reals^n$ we write $Dg (\theta)$ as its Jacobian evaluated at $\theta$.
We say $Dg$ has uniformly lower and upper bounded singular values if there is a constant upper and lower bound to the singular values of $Dg(\theta)$ for all  $\theta \in \reals^d$.
 % We use $\Delta_n =\{x \in \reals^n: \sum_i x^i =1, x^i \geq 0\}$ to denote the $(n-1)$-dimensional simplex.
 


% The surrogate loss approach was originally studied in  supervised learning~\citep{pmlr-v119-johnson20b} and RL ~\citep{vaswani2021general} 
 \paragraph{Surrogate Loss Background.} In the scalar minimization case such as supervised learning, a non-convex loss function of the form $f(g(\theta))$ is minimized where the non-convexity is due to the model parametrization as represented by the model predictions: $g(\theta)= \left[g^1(\theta), \cdots, g^n(\theta)\right]^\top \in \reals^n$.  In supervised learning, each prediction is $z^i = g^i(\theta) = h(x_i, \theta)$, for some feature vector $x_i$ and fixed model architecture $h$.
Despite the non-convexity in parameter space $\theta \in \reals^d$ due to $g$, the loss function is often convex and smooth with respect to the closure of the predictions $\mathcal{Z} = \closure{\{g(\theta) : \theta \in \reals^d\}}$.\footnote{We consider the closure of the predictions to include cases such as softmax paramerizations, where the predictions only lie within the relative interior of the simplex but its closure includes the whole simplex.}
The optimization problem can then be reframed as a constrained optimization problem,
\begin{align}\label{eq:scalar-min}
       \min_\theta f(g(\theta)) = \underset{z\in\mathcal{Z}}{\min} f(z).
\end{align}
If $\mathcal{Z}$ is convex then projected gradient descent $z_{t+1} = \Pi(z_t - \eta \nabla f(z_t))$ is guaranteed to converge~\citep{beck2017first}.
However, the projection $\Pi$ is expensive since it is with respect to the set $\mathcal{Z} \subseteq \reals^n$.
In general, the model-dependent constraint $\mathcal{Z}$ may not be convex. Yet, this assumption is essential since we require convergence of the projected gradient method; we leave relaxing this assumption for future work.  
Nevertheless, this assumption is satisfied in two important extreme cases, when a model is linear or with large capacity neural networks that can interpolate any dataset e.g. $\mathcal{Z}= \reals^n$~\citep{zhang2017understanding}.

Beyond supervised learning, similar hidden structure exists in losses within machine learning such as: generative models and min max optimization~\citep{gidel2021limited}, robust reinforcement learning (RL) \citep{pinto2017robust}, and minimizing projected Bellman error~\citep{bertsekas2009projected}.
However, in these applications the problem cannot be written as minimizing a loss and we must instead consider the VI problem \eqref{eq:vi}.
For example, in the min-max case, the min and max players' strategies may be given by two separate networks $h^1(\theta^1), h^2(\theta^2)$, respectively,
with the following objective:
\begin{align}\label{eq:min-max}
    \min_{\theta^1} \max_{\theta^2}f(h^1(\theta^1), h^2(\theta^2)),
\end{align}
where $f$ is convex-concave. 
Similar to \eqref{eq:scalar-min} we can rewrite the problem in parameter space as a constrained problem with respect to model predictions but instead within the VI \eqref{eq:vi};
where $\theta = (\theta^1, \theta^2)$,  and $g(\theta) = (h^1(\theta^1), h^2(\theta^2))$, with operator $F(z) = F(g(\theta)) = [\nabla_{z^1} f(z^1, z^2), -\nabla_{z^2} f(z^1, z^2)]^\top$ where $z^1= h^1(\theta^1)$ and $z^2= h^2(\theta^2)$.
The solution with respect to model outputs is then equivalent to solving the constrained VI \eqref{eq:vi}.
If $F$ is well-conditioned (e.g. Lipschitz and strongly monotone) and $\mathcal{Z}$ is closed and convex, then the projected gradient method $z_{t+1}=\Pi(z_t -\eta F(z_t))$ converges to a solution $z_\ast$ with an appropriate stepsize $\eta$ \citep{facchinei2003finite}.

To solve problems of the form \eqref{eq:vi} and take advantage of the structure given by $g$ and $F$, we extend the idea of surrogate losses~\citep{pmlr-v119-johnson20b,vaswani2021general}, where parameters are selected by
descending a sequence of surrogates $\ell_t(\theta)$ that approximate the exact projected gradient method.
More precisely, at iteration $t$, $\theta_{t+1}$ is selected by descending the surrogate loss:
\begin{align}\label{eq:surr}
    \ell_t(\theta) &= \frac{1}{2}\big\|g(\theta) - [g(\theta_t) - \eta F(g(\theta_t))] \big\|^2.% \\
    % &= \frac{\eta}{2}\|F(g(\theta_t))\|^2 + \eta\inner{F(g(\theta_t))}{g(\theta)-g(\theta_t)}+ \frac{1}{2}\|g(\theta)-g(\theta_t)\|^2.
\end{align}
Importantly, this loss and its gradient are easily constructed via automatic differentiation packages as a non-linear squared error loss. 
Minimizing the loss exactly gives  $z_{t+1} = g(\theta_{t+1}) =  \Pi(z_t - \eta F(z_t))$, the exact projected gradient step, and guarantees convergence of $\{z_t = g(\theta_t)\}_{t\in\nnumbers}$ to $z_\ast$.

\vspace{-0.5em}
\paragraph{Surrogate losses in scalar minimization.}
The surrogate losses proposed by~\citet{pmlr-v119-johnson20b} and~\citet{vaswani2021general}, apply to supervised learning and RL respectively.\footnote{Note that  \citet{pmlr-v119-johnson20b} and \citet{vaswani2021general} more generally use Bregman divergences.}
They did not study the VI case, nor do they exploit any convexity properties in the scalar minimization case.
\citet{lavington2023target} also study the scalar minimization case and provide convergence to a neighbourhood of a global minimum of \eqref{eq:scalar-min} and allow for stochasticity.
The neighbourhood of convergence depends on an upperbound on the errors $\epsilon_t = \|z_{t+1}-z_t^\ast\|$, where $z_t^\ast$ is an exact projected gradient step.
Therefore, the neighbourhood of convergence scales with the worst error $\epsilon_t$ across the trajectory $\{z_t = g(\theta_t)\}_{t\in \nnumbers}$.
Shrinking the neighbourhood necessitates a double loop algorithm that might spend too much time optimizing the surrogate.
Similar results can also be found within the analysis of quasi-Fej\'er monotone sequences \citep{FRANCI2022161}. 
For example, if $z_t^\ast = T(z_t)$  where $T$ is a contraction, then $\epsilon_t \to 0$ is sufficient to guarantee convergence (see Proposition \ref{thm:contraction-error-limit}).

In contrast to the existing surrogate loss approaches, we propose a simple $\alpha$-descent condition on $\ell_t$ (Definition \ref{eq:alpha}) that does not require all errors to be bounded or summable apriori.
This condition allows for convergence without fully minimizing $\ell_t$ and better models implementations in practice where a fixed number of gradient descent steps are used for each $\ell_t$.  

\begin{definition}[$\alpha$-descent]\label{eq:alpha}Let $\ell_t$ be the surrogate defined in~\eqref{eq:surr} and $\ell_t^\ast = \inf_{\theta \in \reals^d} \ell_t(\theta)$. 
The trajectory $\{\theta_t\}_{t\in \nnumbers}$ satisfies the $\alpha$-descent condition if at each step $t$ the following holds
\begin{align}\label{eq:alpha-ineq}
    \ell_t(\theta_{t+1}) - \ell_t^\ast \leq \alpha^2 \left(\ell_t(\theta_{t}) - \ell_t^\ast \right), \quad \alpha  \in [0, 1).
\end{align}

\end{definition}

Given the $\alpha$-descent condition we can define a general purpose algorithm (Algorithm \ref{alg:surr}).
With any black-box optimizer update $\mathcal{A}$ and a double-loop structure, we can construct a trajectory that satisfies the condition so long as $\mathcal{A}$ can effectively descend the loss $\ell_t$.
In general $\ell_t^\ast$ may not be zero and so this condition cannot be verified directly, however, this condition can often be met via first-order methods for a fixed number of steps or can be approximated with $\ell_t^\ast =0$. 

From a theoretical perspective, we show that the $\alpha$-descent condition guarantees the following:
\begin{align}\label{eq:bias}
    \|z_{t+1} - z_t^\ast \| \leq \alpha \eta \|F(z_t)-F(z_\ast)\|, 
\end{align}
for a proof see Lemma \ref{thm:bias}.
As expected, the errors are controlled by $\alpha$ and the stepsize $\eta$ in the surrogate. However, the error also scales with the distance to the solution if $F$ is Lipschitz.
% , allowing for a more adaptive control of the error.


In the unconstrained scalar minimization case, inequality \eqref{eq:bias} is similar to the relative error condition $\|p_t - \nabla f(z_t)\| \leq \alpha\|\nabla f(z_t)\|$ used in biased gradient descent, where $p_t$ is the biased direction $(z_{t+1}= z_t -p_t)$~\citep{ajalloeian2020convergence,drusvyatskiy2023stochastic}.
% If $z_t^\ast = z_t-\eta F(z_t)$ (i.e. $\mathcal{Z}$ is unconstrained) then the above inequality guarantees a positive inner product with a cosine angle of at least $\sqrt{1-\alpha^2}$ between the directions $p_t$ and $\eta F_t$ where $p_t$ is the actual direction of change, $z_{t+1}= z_t -p_t$~\citep[Appendix E]{drusvyatskiy2023stochastic}.
Similar to biased gradient descent, our assumption with any $\alpha <1$ is sufficient to guarantee convergence in scalar minimization with any $L$-smooth  non-convex loss $f$~(see Proposition \ref{thm:alpha-scalar}).
% For any $L$-smooth  non-convex loss $f(z)$, the $\alpha$-descent  condition guarantees  $f(z_{t+1}) \leq f(z_t) - L(1-\alpha^2)(\ell_t(\theta_t) - \ell_t^\ast)$(see Proposition \ref{thm:alpha-scalar}), if $f$ is bounded below then we have $z_t \to z_t^\ast$.
However, for VIs, $\alpha <1$ can still lead to divergence due to possible rotations in the dynamics of $z_t^\ast$ (see Proposition~\ref{prop:counter}).

Within the context of solving VI problems,  \citet{solodov1999hybrid} proposed a similar condition to approximate the proximal point algorithm where $z_t^\ast = z_t - \eta F(z_t^\ast)$.\footnote{They assume $\|z_{t+1} -(z_t-\eta F(z_{t+1}))\| \leq \alpha \|z_{t+1}-z_t\|$.}
However, verifying their condition would  induce an inner loop with multiple evaluations of $F$, which can be expensive in many applications such as reinforcement learning~\citep{vaswani2021general}.
They also show that divergence is still possible under their condition unless an extra-gradient like step is performed.
Unfortunately, such a step is impractical in our setting as it requires minimizing a surrogate loss exactly. 

\paragraph{Hidden monotone problems and preconditioning methods.}
In scalar minimization problems, hidden monotone structure has been studied under hidden convexity~\citep{JMLR:francis,xia2020survey, chancelier2021conditional, fatkhullin2023stochastic}.
In zero-sum games with hidden monotonicity, \citet{gidel2021limited} study the existence of equilibria and establish approximate min-max theorems but do not propose an algorithm to take advantage of the hidden structure.
For games that admit a hidden strictly convex-concave structure, \citet{vlatakis2021solving} prove global convergence of continuous time gradient descent-ascent.
Similarly, \citet{mladenovic2022generalized} propose natural hidden gradient dynamics (NHGD) with continuous time global convergence guarantees in hidden convex-concave games.
A more general and descritized version of NHGD was studied by \citet{sakos2024exploiting}, the preconditioned hidden gradient descent method (PHGD), to solve VIs of the form \eqref{eq:vi}. One step of PHGD corresponds to the update:
\begin{align}\label{eq:phgd}
    \theta_{t+1} = \theta_t - \eta(Dg(\theta_t)^\top Dg(\theta_t))^{\dag}Dg(\theta_t)^\top F(z_t).
\end{align}
PHGD and stochastic variants were also studied in the linear case by \citet{bertsekas2009projected}.
PHGD can be viewed as a discretization of a continuous flow that guarantees $z_{t+1} = z_t-\eta F(z_t)$ as $\eta \downarrow 0$ \citep{sakos2024exploiting}. 
Interestingly, in Section \ref{sec:non-linear} we show that PHGD is also equivalent to taking one step of the Gauss-Newton method (GN)~\citep{bjorck1996numerical} on the surrogate loss.


% Despite the nice properties of PHGD, it inherits some of the instabilities of GN.
% These instabilities can be controlled by using alternatives to GN such as damped Gauss-Newton (DGN) or the Levenbergh-Marquardt (LM) method and its variants ~\citep{bjorck1996numerical, nocedal1999numerical, mishchenko-reg-newton}.
% However, these methods require computing the pseudo-inverse of a potentially large matrix and therefore are not scalable if the prediction space or parameter space is large. 


\vspace{-0.5em}
\section{Convergence analysis under $\alpha$-descent on surrogate losses}

In this section we provide analysis in both the deterministic and stochastic settings.
In both settings we make the following assumption on the hidden structure in the VI \eqref{eq:vi}.
\begin{assumption}\label{assumption:vi}
    In the VI \eqref{eq:vi}, $\mathcal{Z}$ is convex. There exists a solution within the relative interior, $z_\ast \in \ri \mathcal{Z}$.
    $F$ is both $L$-Lipschitz $\|F(x)-F(y)\|\leq L \|x-y\|$ and $\mu$-strongly monotone $\inner{F(x)-F(y)}{x-y} \geq \mu \|x-y\|^2$ for any $x,y \in \mathcal{Z}$ and some $\mu > 0$.
\end{assumption}
This assumption is commonly used in the VI literature to establish linear convergence of the projected gradient method $z_t = \Pi(z_t- \eta F(z_t))$~\citep{facchinei2003finite}.
In the scalar minimization case this assumptions implies $F(z) = \nabla f(z)$ for a scalar loss function $f$ that is $\mu$-strongly-convex and $L$-smooth.
In min-max optimization with hidden structure, such as \eqref{eq:min-max}, Assumption \eqref{assumption:vi} is satisfied if $f(z^1, z^2)$ is strongly-convex-concave and smooth (each players' gradient is Lipschitz with respect to all players, e.g. see ~\citet{bubeck2015convex}). It corresponds to many practical cases as the losses used in machine learning applications are often strongly convex \emph{with respect to the model predictions $z$}. note that $\mathcal{Z}$ is convex for many classes of models used in practice, e.g., linear models or models that can interpolate any noisy labels on the train set ($\mathcal{Z}= [0,1]^n$). The latter has been reported with large enough neural networks~\citep{zhang2017understanding} and kernels~\citep{belkin2019does}.


% In the deterministic setting, we provide linear convergence guarantees of Algorithm \ref{alg:surr} using the $\alpha$-descent condition (Definition \eqref{eq:alpha}) if $\alpha$ is small enough. 
% We show that a small enough $\alpha$ 

% guarantees in the deterministic setting with Algorithm \ref{alg:surr} and the $\alpha$-descent condition \eqref{eq:alpha}.
% In the stochastic case we use and the expected $\alpha$-descent condition
% Under the $\alpha$-descent condition \eqref{eq:alpha} we provide convergence guarantees to the solution $z_\ast$ of the VI \eqref{eq:vi}.
% Throughout this section, we use the following assumption:



\subsection{Convergence and divergence in the deterministic case}

\paragraph{Convergence for sufficiently small $\alpha$.} Below we provide a linear convergence result if $\alpha$ is small enough by controlling the stepsize $\eta$ in the surrogate loss $\ell_t$.
Note that the convergence is with respect to $z_t = g(\theta_t)$ and is linear in the number of outer loop iterations of Algorithm \ref{alg:surr}.
% If $\alpha =0$ we recover the standard rate of the projected gradient method.
\begin{theorem}\label{thm:det}
    Let Assumption \ref{assumption:vi} hold and let $\{z_t = g(\theta_t)\}_{t\in \nnumbers}$ be  the iterates produced by Algorithm~\ref{alg:surr}. If $\alpha$ and $\eta$ are picked  such that
$
\rho:= 1-2\eta(\mu-\alpha L) + (1+\alpha^2)\eta^2 L^2 <1 
$ then, $z_t$ converge linearly to the solution $z_\ast$ at the following linear rate:
\begin{equation}
    \|z_{t+1} -z_*\|^2 \leq \rho^t \|z_1-z_*\|^2.
\end{equation}Particularly, if $\alpha< \frac{\mu}{L}$ and $\eta < \frac{2 (\mu-\alpha L)}{(1+\alpha^2)L^2}$ then $\rho <1$ and if
$\alpha \leq \frac{\mu}{2L}$ and $\eta = \frac{2\mu}{5L^2}$ then $\rho \leq 1 - \frac{\mu^2}{5L^2}$.
\end{theorem}

Note that, to obtain a practical convergence rate (e.g., bounding the number of gradient computations), we need to bound the number of inner-loop steps in Algorithm~\ref{alg:surr} required to obtain the condition~\eqref{eq:alpha-ineq}. In general, we cannot provide any global guarantee as it is well established that, in general, finding a global minima of a smooth non-convex optimization function like~\eqref{eq:surr} can be intractable~\citep{murty1985some,nemirovskij1983problem}. However, many classes of overparametrized models are known to be able to \emph{interpolate} random labels (e.g., neural networks~\citep{zhang2017understanding}, kernels~\citep{belkin2019does}, or boosting~\citep{bartlett1998boosting}) which is strong evidence that, in practice, the condition~\eqref{eq:alpha-ineq} can be obtained with a few gradient steps. The latter is also supported by our experiments in Section~\ref{sec:experiments}.
% Alternatively, for a fixed stepsize $\eit's ota$, $\alpha$ can be set small enough that guarantees convergence.
\vspace{-1em}
\paragraph{Divergence with $\alpha < 1$.}
% \textcolor{red}{Comment more why it is interesting. It is a counterexample that $\alpha<1$ is not sufficient}
In \emph{non-convex} smooth scalar minimization, $\alpha < 1$ is sufficient for convergence~(Proposition \ref{thm:alpha-scalar}). 
However, despite our strong monotonicity assumption we show that $\alpha < 1$ can still give divergence in the VI setting. 
This shows that $\alpha$ being small enough is not only sufficient but is \textit{necessary} for convergence.
Our example demonstrates the additional challenges of using a surrogate loss in the VI setting.
Our construction uses the min-max problem  $\min_x \max_y \frac{1}{2}x^2 + xy -\frac{1}{2}y^2$ and showing that gradient descent-ascent on the loss $f(x,y) = xy$ satisfies the $\alpha$-descent condition with $\alpha =\nicefrac{1}{\sqrt{2}}$ but diverges for any $\eta$.

\begin{proposition}\label{prop:counter}
There exists an $L$-Lipschitz and $\mu$-strongly monotone $F$, and a sequence of iterates $\{z_t\}_{t \in \nnumbers}$ verifying the alpha descent condition with  $\alpha < 1$  such that  $z_t$ diverges for any $\eta$.
\end{proposition}


\subsection{Unconstrained Stochastic Case}\label{sec:stochastic}
For the stochastic case we assume $\mathcal{Z}=\reals^n$, the predictions can represent all of $\reals^n$.
Although this assumption is strong it can be satisfied for large capacity neural networks that can interpolate any dataset. 
The stochastic case is more challenging as our setup corresponds to solving VIs with bias, for which we show can diverge even in the deterministic case (Proposition \ref{prop:counter}).

Since $\mathcal{Z}=\reals^n$, no projection is needed for $z_t^\ast$, and the minimum of the surrogate $\ell_t$ is the gradient step $z_t^\ast = z_t - \eta F(z_t)$.
We also assume the standard setup, the possibility to generate independent and identically distributed realizations $(\xi_1, \xi_2, \cdots)$ of a random variable $\xi$ such that $F_{\xi}(z)$ is an unbiased estimator of $F(z)$.
For the noise, we define $\sigma^2 = \expect[\xi]{\|F_\xi(z_\ast)\|^2}$ and use the expected co-coercive assumption from ~\citet{loizou2021stochastic}:
 \begin{assumption}[expected co-coercivity\footnote{Note that this is not exactly expected co-coercivity but implied by it~\citep[Lemma 3.4]{loizou2021stochastic}.}]\label{eq:exp-co}   $\expect[\xi]{\|F_\xi(z)\|^2} \leq 2L\inner{F(z)}{z-z_\ast}+2\sigma^2,$ for all $z \in \mathcal{Z}$. 
 \end{assumption}
Without access to the true operator $F$ in the stochastic case, a different $\alpha$-descent condition is required using the stochastic estimate $F_{\xi_t}(z_t)$.
Instead of $\ell_t$ \eqref{eq:surr}, we will use the loss $\Tilde{\ell}_t$, an approximation of a projected stochastic gradient step, $\Tilde{\ell}_t(\theta) = \frac{1}{2}\|g(\theta)-(g(\theta_t)-\eta F_{\xi_t}(z_t))\|^2$, or equivalently
\begin{align}
    \Tilde{\ell}_t(\theta) = \frac{\eta^2}{2}\|F_{\xi_t}(z_t)\|^2 + \eta\inner{F_{\xi_t}(z_t)}{g(\theta)-g(\theta_t)}+\frac{1}{2}\|g(\theta)-g(\theta_t)\|^2.\label{eq:surr-stoch}
\end{align}
In practice, minimizing $\Tilde{\ell}_t$ exactly is impractical if $n$ is very large due to the sum of squares $\|g(\theta)-g(\theta_t)\|^2$ in \eqref{eq:surr-stoch}. Therefore we assume it can be minimized on expectation with the following assumption. 
% $\alpha$-expected descent condition (Definition \ref{eq:exp-bias}).

\begin{definition}[$\alpha$-expected descent]\label{eq:exp-bias}
The trajectory $\{\theta_t\}_{t\in \nnumbers}$ satisfies the $\alpha$-expected descent condition if at for each $t$ the following holds:
given $F_{\xi_t}(z_t)$,
the parameter $\theta_{t+1}$ is generated such that 
\[
\expect{\Tilde{\ell}_t(\theta_{t+1})-\Tilde{\ell}_t^\ast} \leq \alpha^2(\Tilde{\ell}_t(\theta_{t})-\Tilde{\ell}_t^\ast).
\]
\end{definition}
% Similar to~\eqref{eq:bias} we have the bound $\expect{\|z_{t+1}-z_t^\ast\|^2} \leq \alpha^2\eta^2\|F_{\xi_t}(z_t)\|^2$, and we can prove linear convergence on expectation to a neighbourhood. 
Similar to the deterministic case, we can prove linear convergence to a neighbourhood.
\begin{theorem}\label{thm:stochastic}
Let $\mathcal{Z} = \reals^n$, and Assumption~(\ref{eq:exp-co}) hold.  
If $F_{\xi}(x)$  is $L$-Lipschitz and $\eta \leq \frac{1}{2(1+c)L}$ where  $c \geq  2(1+\alpha^2)$  then any trajectory $\{\theta_t\}_{t\in \nnumbers}$ 
satisfying the $\alpha$-expected descent condition guarantees:
\[
\expect{\tfrac{1}{2}\|z_{t+1}-z_\ast\|^2} \leq \expect{\tfrac{1}{2}\|z_{t}-z_\ast\|^2}\left(1-\eta\mu + \alpha^2\right) +\eta^2(1+c)\sigma^2. 
\]
\end{theorem}
 If we take $c=4$ and $\eta\leq \nicefrac{1}{10 L}$ then we have convergence to a neighbourhood if $\alpha < \sqrt{\eta \mu}$.
 Theorem \ref{thm:stochastic} provides a generalization of~\citet{loizou2021stochastic}[Theorem 4.1] where $z_{t+1}$ follows a random biased direction $p_t$ that on expectation guarantees $\expect{\|p_t -\eta F_{\xi_t}(z_t)\|^2} \leq \alpha^2 \eta^2 \|F_{\xi_t}(z_t)\|^2$.
 In comparison to the stochastic results of \citet{lavington2023target}, they avoid the assumption $\mathcal{Z} = \reals^n$ by using a different surrogate. However, their result only applies to scalar minimization.

 
 \section{A Nonlinear Least Squares Perspective}\label{sec:non-linear}
The surrogate loss perspective and our $\alpha$-descent condition allows for convergence so long as the surrogate losses $\{\ell_t\}_{t\in \nnumbers}$ are sufficiently minimized.
One approach to minimizing $\ell_t$ is to view it as the following  non-linear least-squares problem 
\begin{align}
\min_\theta f(\theta) = \min_\theta \frac{1}{2}\|r(\theta)\|^2,
\end{align}
with a residual function $r: \reals^d \to \reals^n$,
where $\ell_t(\theta)= f(\theta)$ if $r(\theta) = g(\theta)-g(\theta_t)+\eta F(g(\theta_t))$.
Due to the specific structure of $f$ we can consider  specialized methods such as Gauss-Newton (GN), Damped Gauss-Newton (DGN), and Levenbergh-Marquardt (LM)~\citep{bjorck1996numerical,nocedal1999numerical}.
These methods can be viewed as quasi-Newton methods that use a linear approximation of $r$, $r(\theta) \approx r(\theta_t) + Dr(\theta_t)(\theta -\theta_t)$. 

The GN method is defined by the update rule $\theta_{t+1} = \theta_t - (Dr(\theta_t)^\top Dr(\theta_t))^\dag Dr(\theta_t)^\top r(\theta_t)$. 
GN inherits the same local quadratic convergence properties as Newton's method
when the Hessian at the minimum $\nabla^2 f(\theta_\ast) \approx Dr(\theta_\ast)^\top Dr(\theta_\ast)$. 
However, GN is known to struggle with highly non-linear problems, those with large residuals, or if $Dr(\theta_t)$ is nearly rank-deficient~\citep{bjorck1996numerical}.
Fortunately, the GN direction is a descent direction of $f$, the DGN method takes steps in the GN direction with a stepsize parameter $\eta_{GN}$ and converges for a sufficiently small stepsize or with line search~\citep{bjorck1996numerical}.
In cases where $Dr(\theta_t)$ is nearly rank deficient, the LM method can be used. 

To minimize the surrogate we can consider taking multiple steps of gradient descent (Surr-GD), DGN or LM. 
Denoting $\theta_t^s$ as $s$\textsuperscript{th} intermediate step between $\theta_{t+1}$ and $\theta_t$ we have the following:
\begin{align}
        \tag{Surr-GD} \theta_t^{s+1} &= \theta_t^s - \eta_{GD} Dg(\theta_t^s)^\top (g(\theta_t^s)-g(\theta_t)+\eta F(g(\theta_t))\\
    \tag{DGN}\label{eq:DGN} \theta_t^{s+1} &= \theta_t^s - \eta_{GN} (Dg(\theta_t^s)^\top Dg(\theta_t^s))^{\dag} Dg(\theta_t^s)^\top (g(\theta_t^s)-g(\theta_t)+\eta F(g(\theta_t))\\
    \tag{LM} \theta_t^{s+1} &= \theta_t^s - (Dg(\theta_t^s)^\top Dg(\theta_t^s) +\lambda \id)^{-1} Dg(\theta_t^s)^\top (g(\theta_t^s)-g(\theta_t)+\eta F(g(\theta_t)).
\end{align}
Note that we used the fact that $Dr(\theta) = Dg(\theta)$, and if $\eta_{GN}=1$ then DGN is the same as GN.
Also, note that one step of GN recovers exactly the PHGD method proposed by \citet{sakos2024exploiting}.

\subsection{Favourable conditions for gradient descent}\label{sec:gd-gn}

Several conditions allow for fast linear convergence of gradient descent for $\ell_t(\theta) - \ell_t^\ast$; see for example ~\citet{guille-escuret21a}.
Under linear convergence, $s$-steps of gradient descent guarantees $\alpha = \rho^s$ for some $\rho \in [0,1)$, therefore any target value of $\alpha$ is achievable in a finite number of steps that depends only logarithmically in $\alpha$.
In this section, we show how assumptions used in \citet{sakos2024exploiting} imply two common assumptions for fast convergence of GD: the Polyak-\L{}ojasiewicz (PL) condition ~\citep[Definition  \ref{def:pl}]{polyak1964gradient,lojasiewicz1963propriete} and $L$-smoothness.

 The four assumptions made by \citet{sakos2024exploiting} to study the behaviour of PHGD are: (1-2) $Dg^\top$ has both uniformly lower bounded and upper bounded singular values, (3) each component function $g_i$ is $\beta$-smooth, (4) and that $\mathcal{Z}$ is bounded.\footnote{
 The set $\mathcal{Z}$ (note that $\mathcal{Z}$ is the ``set of latent variables" denoted as $\mathcal{X}$ in  \citet{sakos2024exploiting}) is assumed to be bounded in their Lemma 4 ``template inequality". More precisely, the existence of a constant $D = diam(\mathcal{X})$ (i.e. $\mathcal{X}$ is bounded) is used in equation B.20 in Appendix B. Lemma 4 is then used in Theorems 1-4. 
 % Although $\mathcal{Z}$ being bounded is not explicitly stated as an assumption it is used within the their main Lemma.
 }
In Proposition \ref{thm:pl} we show that globally lower bounding the singular values of $Dg^\top$ guarantees that the composition of a PL function $f$ with $g$ is still PL.
\begin{proposition}\label{thm:pl}
    Assume $f$ satisfies the $\mu$-PL condition, and let $\sigma_{\min}$ be a lower bound on the singular values of $Dg(\theta)^\top$. Then, $f \circ g$ is $\mu \sigma^2_{\min}$-PL.
\end{proposition}
This implies that the surrogate loss $\ell_t$ is PL since $\ell_t = f_t \circ g$ where $f_t(z) = \nicefrac{1}{2}\|z-v_t\|^2$ is $1$-PL.
Similarly, $\ell_t$ can be shown to be $L$-smooth if $Dg^\top$ has uniformly upper bounded singular values, each $g_i$ is $\beta$-smooth, and $\mathcal{Z}$ is bounded (see Lemma \ref{thm:smooth-surr}).
However, these 4 assumptions cannot hold all at once. In Prop.~\ref{thm:unbounded}, we show that if the first three assumptions (1-3) hold, then $\mathcal{Z}$ is unbounded.
\begin{proposition}\label{thm:unbounded}
    If $g: \reals^d \to \reals^n$ is differentiable where $g_i$ is $\beta$-smooth and $Dg(\theta)^\top$ has globally lower and upper bounded singular values, then $\{g(\theta): \theta\in \reals^d\}$ is unbounded.
\end{proposition}

This contradiction suggests that the lower bounded singular value assumption is strong.
Indeed, it forces the dimension of the parameter space $d$ to be larger than the dimension of the prediction space $n$, and if $g$ is linear then it must be surjective $\mathcal{Z} = \reals^n$.
In the case $n=d$, it enforces $Dg$ to be invertible everywhere, which is violated in important cases such as softmax. 

\section{Experiments}
\label{sec:experiments}

We consider min-max optimization and policy evaluation to demonstrate the behaviour and convergence of surrogate based methods with hidden monotone structure.
For min-max optimization we compare different approaches from Section \ref{sec:non-linear} on two domains from \citet{sakos2024exploiting}: hidden matching pennies and hidden rock-paper scissors.
For RL policy evaluation we consider the problem of minimizing projected Bellman error (PBE) with linear and non-linear approximation.  

\subsection{Min-max Experiments}
\label{sec:experiments_minmax}

\begin{figure}
    \centering
        \vspace{-13mm}
    \includegraphics[width=\linewidth]{figures/figure_1_pennies_crop.pdf}
    \caption{
    \small
    Convergence of various algorithms from Section \ref{sec:non-linear} on the hidden matching pennies game. PHGD and GDA as presented in \citet{sakos2024exploiting} are compared against GN, DGN, LM, and GD. (left) Linear convergence to the equilbrium is observed for several methods with LM and GD outperforming the rest. (middle) Trajectories for some methods  are plotted in both the parameter and prediction space. (right) The loss ratio $\ell_t(\theta_{t+1})/\ell_t(\theta_t)$ is illustrated for the considered methods.}
    \label{fig:pennies}
    \vspace{-5mm}
\end{figure}

We compare four different approaches from Section ~\ref{sec:non-linear}: GN, DGN, LM, and Surr-GD.
Taking only one inner step for GN and Surr-GD gives PHGD and gradient descent-ascent (GDA), respectively.

\begin{wrapfigure}[13]{R}{0.55\textwidth}
\vspace{-0.5em}
\centering
\includegraphics[width=0.9\linewidth]{figures/figure_2_rps_crop.pdf}
\vspace{-0.75em}
\caption{
\small
Convergence in the hidden rps game.}
\label{fig:rps}
\end{wrapfigure}

\paragraph{Hidden Matching Pennies.} The hidden matching pennies game is a zero-sum game of the form \eqref{eq:min-max}, where each player has the parameterization $h^i(\theta) = \sigmoid(\alpha_2^i\celu(\alpha^i_1 \theta))$.
The convex-concave objective $f$ is given by $f(z^1,z^2)= -(2z^1-1)(2z^2-1)+\frac{0.75}{2}\left((z^1-\nicefrac{1}{2})^2 + (z^2-\nicefrac{1}{2})^2\right)$.
The parameters $\alpha_j^i$ are chosen to approximately replicate the trajectory of PHGD presented in \citet[Figure 4]{sakos2024exploiting} (see Appendix \ref{sec:appendix-pennies} for more details).
In Figure \ref{fig:pennies}, we observe that PHGD converges linearly in the squared distance to the equilibrium, but performs poorly if multiple inner steps are taken (GN with 5 inner steps). 
If $\eta$ is increased by an order of magnitude, PHGD is observed to diverge; however, convergence is possible via DGN with the same $\eta$, with multiple iterations and an appropriate choice of $\eta_{DGN}$.
In contrast to GN, LM is more stable with a larger $\eta$ and converges faster than PHGD/GN. 
Finally we tested GD for a different number of inner steps with $\eta=0.1$.
Convergence is observed for GDA albeit slow.
The benefit of multiple steps is clear, with 10 inner steps outperforming PHGD and only surpassed by LM.
Although more inner steps increases the computational cost, it is marginal when compared to evaluating $F$ (see Fig.~\ref{fig:pennies-time}).
Interestingly, Fig.~\ref{fig:pennies} (right) shows that spending more compute to minimize the surrogate at each iteration does not necessarily translate to faster overall convergence with respect to the outer loop (left).
GD with 10 inner steps has a larger loss ratio than GD with 100 steps but converges faster to the equilibrium.

\paragraph{Hidden rock-paper-scissors.} In the hidden rps game each player's mixed strategy in rock-paper-scissors is parameterized. 
Where player $i$'s strategy $z^i$ is given by the function $h^i(\theta) = \softmax(A_2^i \celu(A_1^i \theta^i))$, with  $\theta^i \in \reals^5$ and randomly initialized matrices: $A_1^i \in \reals^{4 \times 5}$, and $A_2^i \in \reals^{3 \times 4}$.
Figure \ref{fig:rps} demonstrates the behaviour of various algorithms for a fixed initialization of $\theta = (\theta^1, \theta^2)$ and the matrices $A^i_j$. 
We observe that PHGD and LM with one inner step achieve linear convergence while GDA performs poorly with an unstable behaviour.
Like in the hidden matching pennies game, increasing the number of inner steps for GD improves stability and performance, with the best performance not necessarily corresponding to the methods with the lowest loss ratio.
Both LM and GD degrade in performance if too many inner steps are taken, with one and 10 inner steps outperforming 5 and 100 steps for LM and GD respectively.


\subsection{Minimizing Projected Bellman error}
\label{sec:experiments_PBE}

For our RL experiments we consider the policy evaluation problem of approximating a value vector $v_{\pi} \in \reals^n$, representing the expected discounted return over $n$ states for a given policy $\pi$ and Markov decision process (MDP).
To this end, we consider a common approach to policy evaluation: minimizing the projected Bellman error, which is a fixed point problem associated with temporal difference learning (TD)~\citep{sutton1988learning}. %, a common approach to policy evaluation.
Despite the widespread practical success of TD and its variants, the analysis and behaviour of TD has proven challenging since it cannot be modeled as a scalar minimizing problem and is known not to follow the gradient of any objective~\citep{baird1995residual, antos2008learning}.
\citet{bertsekas2009projected} showed that minimizing PBE is equivalent to solving a smooth and strongly monotone VI problem of the form \eqref{eq:vi} where $F(z) = \Xi(z-T_{\pi}(z))$~\citep{bertsekas2009projected}. $T_{\pi} : \reals^n \to \reals^n$ is the linear Bellman policy evaluation operator defined over all $n$ states $(s^1, \cdots, s^n)$ in a given MDP and is defined as $T_{\pi}(z) = r_{\pi} + \gamma P_{\pi}z$,
where $r_{\pi}$ is the associated expected reward vector at each state and $P_{\pi}$ is the state to state probability transition matrix as given by the MDP and $\pi$.
$\Xi$ is the diagonal matrix with entries corresponding to the stationary distribution $\xi \in \reals^n$ of states according to $\pi$.
The constraint $\mathcal{Z}$ in this case is the set of all representable value functions. 
With respect to $\theta$, solving the VI associated with minimizing PBE is equivalent to finding the fixed point:
\begin{align}\label{eq:PBE}
    \theta_\ast \in \underset{\theta}{\arg\min} \left[ BE(\theta, \theta_\ast) =  \frac{1}{2}\sum_{i=1}^n \xi^i(v_{\theta}(s^i) - (r_\pi(s^i)- \gamma\expect[s'\sim P_\pi]{v_{\theta_\ast}(s')|s^i})^2\right],
\end{align}

where $v_{\theta}$ is the predicted value vector given parameters $\theta$. 
% Notice that $BE(\theta, \theta')$ is the surrogate loss $\nicefrac{1}{2}\|v_{\theta} - T_{\pi}(v_{\theta'})\|^2_{\Xi}$.
See Appendix \ref{sec:appendix-PBE} for more details.

\subsubsection{Linear Approximation}
\begin{wrapfigure}[15]{R}{0.45\textwidth}
\centering
\vspace{-3em}
\includegraphics[width=\linewidth]{figures/bert_exp_crop.pdf}
\vspace{-2em}
\caption{\small The average approximation error between GD on the surrogate (\ref{eq:surr-stoch-berts}, Surr-GD) and update \eqref{eq:bert-update-stoch} over 10,000 runs in a slow mixing 100-state Markov chain from \citet{bertsekas2009projected} and \citet[Example 3]{Hu_Berts_markov_chains_ex}.
Surr-GD is observed to converge to the exact update \eqref{eq:bert-update-stoch} with faster convergence for more inner steps.}
\label{fig:bertsekas_experiments}
% \vspace{-2cm}
\end{wrapfigure}


For our linear RL experiments, we consider a slow-mixing 100-state Markov chain from \citet{bertsekas2009projected} and \citet{Hu_Berts_markov_chains_ex}.
Using a linear model, $z = g(\theta) = \Phi \theta$ with $\theta \in \reals^d$, $\Phi \in \reals^{n \times d}$ and $d \ll n$.
\citet{bertsekas2009projected} proposed several methods to leverage the hidden structure in $F$.
Similar to \citet{sakos2024exploiting}, the proposed methods are presented via preconditioning schemes that approximate a projected gradient step in the space of representable value functions $\mathcal{Z}$.
\citet{bertsekas2009projected} suggests performing the deterministic and stochastic PHGD-like updates:
\begin{align}
    \theta_{t+1} &= \theta_t - (\Phi^\top\Xi\Phi)^{-1}\Phi^\top F(z_t), \label{eq:bert-update} \\
    \theta_{t+1} &= \theta_t - \hat{D}_t^{-1} (\hat{C}_t\theta_t -\hat{r}_t). \label{eq:bert-update-stoch}
\end{align}

The deterministic update is expensive, requiring the full feature matrix $\Phi$. Hence, \citet{bertsekas2009projected} constructs an online sequence of estimators $\{\hat{D}_t\}_{t\in \nnumbers}$, $\{\hat{C}_t\}_{t\in \nnumbers}$, and $\{\hat{r}_t\}$ from a trajectory and suggests the stochastic update (\ref{eq:bert-update-stoch}). 
Although this method applies more generally, in the context of solving \eqref{eq:PBE} update \eqref{eq:bert-update-stoch} is equivalent to least-squares policy evaluation~\citep{bertsekas1996temporal}.

Equivalently, update \eqref{eq:bert-update} can be seen as one step of the generalized Gauss-Newton method~\citep{ortega2000iterative} on the surrogate $\ell_t(\theta) = \nicefrac{1}{2}\|\Phi \theta - T_{\pi}(\Phi\theta_t)\|^2_{\Xi}$. 
Since the model is linear, the update is the minimizer of the surrogate, and corresponds to an exact projected gradient step in the space of value functions, which is known to be a contraction~\citep{bertsekas2012dynamic}.

% Notice that Equation \ref{eq:bert-update} requires storing the feature matrix $\Phi$ and implementing the operator $T_\pi$.
%Instead of performing the exact deterministic update, \citet{bertsekas2009projected} constructs an online sequence of estimators $\{\hat{D}_t\}_{t\in \nnumbers}$, $\{\hat{C}_t\}_{t\in \nnumbers}$, and $\{\hat{r}_t\}$ from a trajectory and suggests the stochastic update (\ref{eq:bert-update-stoch}).
Similar to the deterministic update, we show that the stochastic version can be shown to be the minimizer of a surrogate loss $\tilde{\ell}_t(\theta)$, whose minimizer is the exact update (\ref{eq:bert-update-stoch}).
The stochastic surrogate $\tilde{\ell}_t$ and its gradient, derived in Proposition \ref{thm:bert-surr}, are
\begin{align} \label{eq:surr-stoch-berts}
    \begin{aligned}
        \tilde{\ell}_t(\theta) = \frac{1}{2}\|\Phi \theta - \hat{T}_{\pi}(z_t)\|^2_{\hat{\Xi}},
    \end{aligned}
    && 
    \begin{aligned}
        \nabla \tilde{\ell}_t(\theta) = \hat{C}_t \theta_t-\hat{r}_t + \hat{D}_t(\theta - \theta_t),
    \end{aligned}
\end{align}
where $\hat{\Xi}$ and $\hat{T}_{\pi}$ correspond to the natural estimators using the empirical distribution (see Appendix ~\ref{sec:appendix-linear} for more details). Instead of minimizing the loss exactly it can be approximated with GD, thus avoiding the matrix inversion $\hat{D}^{-1}_t$, which can make \eqref{eq:bert-update-stoch} intractable. 
Figure \ref{fig:bertsekas_experiments} demonstrates the effectiveness of approximating update \eqref{eq:bert-update-stoch} via GD on the surrogate $\tilde{\ell}_t$~\eqref{eq:surr-stoch-berts}.


% \vspace{-7pt}



% Update \eqref{eq:bert-update} corresponds to the doing an exact projected gradient step in the space of representable value functions, $z_{t+1}=\Pi_{\Xi}(z_t -\Xi^{-1} F(z_t)) = \Pi_{\Xi}(T_{\pi}(z_t))$, which is known to be a contraction~\citep{bertsekas2012dynamic}.  
% Equivalently, $\theta_{t+1}$ is the minimizer of the surrogate $\ell_t(\theta) = \nicefrac{1}{2}\|\Phi \theta - T_{\pi}(\Phi\theta_t)\|^2_{\Xi}$ and can be viewed as performing one step of the generalized Gauss-Newton method~\citep{ortega2000iterative}. 
% If $\Xi$ is the identity then update \eqref{eq:bert-update} reduces to PHGD. 

% Despite the convergence properties of update \eqref{eq:bert-update} it is impractical because it requires storing the feature matrix $\Phi$ and performing an update with the operator $T_{\pi}$.
% Instead of performing the exact update, \citet{bertsekas2009projected} suggests constructing an online sequence of estimators $\{\hat{D}_t\}_{t\in \nnumbers}$, $\{\hat{C}_t\}_{t\in \nnumbers}$, and $\{\hat{r}_t\}$ from a trajectory and proposed the update
% $\theta_{t+1} = \theta_t - \hat{D}_t^{-1} (\hat{C}_t\theta_t -\hat{r}_t)$.
% % \citet{bertsekas2009projected} also proposes scaling with the diagonal approximation $\diag(\hat{D}_t)^{-1}$ and $(\hat{D}_t +\lambda \id)^{-1}$ similar to the LM method.
% Similar to the deterministic update, the stochastic version can be shown to be the minimizer of the surrogate $\tilde{\ell}_t(\theta) = \frac{1}{2}\|\Phi \theta - \hat{T}_{\pi}(z_t)\|^2_{\hat{\Xi}}$ where $\hat{\Xi}$ and $\hat{T}_{\pi}$ corresponds to the natural estimators using the empirical distribution, for full details see Section \ref{sec:appendix-PBE}.

% Similar to Bertsekas' method ~\eqref{eq:bert-update}, gradient descent on the surrogate $\ell_t$ is intractable with a large amount of states but  is tractable on the loss $\tilde{\ell}_t$  and can be updated online with the same estimators
% $\nabla \tilde{\ell}_t(\theta) = \hat{C}_t \theta_t-\hat{r}_t + \hat{D}_t(\theta - \theta_t)$(see Proposition \ref{thm:bert-surr}).
% Instead of minimizing the loss exactly, which requires inverting the matrix $\hat{D}_t$, we can perform GD on the surrogate.




\subsubsection{Non-linear Projected Bellman Error}
\begin{figure}
    \centering
        \vspace{-13mm}
    \includegraphics[width=\linewidth, scale=0.4]{figures/Figure_1.pdf}
    \vspace{-7mm}
    \caption{
    \small
    Comparison of average performance of TD(0) and surrogate methods in minimizing the value prediction error for RL tasks with nonlinear function approximation in Ant (top) and HalfCheetah (bottom) environments, measured by outer loop iterations (left) and wallclock time (right).
    The average value prediction error across 20 runs along with 95\% confidence intervals are computed from a fixed test set.
    The red dashed line represents the lowest value prediction error achieved by any of the algorithms.
    }
    \label{fig:nonlinear}
        \vspace{-7mm}
\end{figure}

For the non-linear setting, we approximate the value function of a fixed reference policy with a two-layer neural network in two Mujoco environments, Ant and Half Cheetah~\citep{todorov2012mujoco}. These environments are characterized by continuous state and action spaces, where the agent must control high-dimensional, physics-based simulations to achieve locomotion tasks. The reference policy is derived by training a policy network with the soft-actor-critic algorithm~\citep{haarnoja2018soft}.
For each environment, we consider two different surrogate-based methods with TD(0) as a special case.
Similar to the linear case, we aim to find the fixed point \eqref{eq:PBE}, and consider the surrogate $\ell_t(\theta) = \nicefrac{1}{2}\|v_{\theta}-T_{\pi}(v_{\theta_t})\|_{\Xi}^2$.
However, due to the large state spaces we must instead consider miniziming the stochastic surrogate $\Tilde{\ell}_t$ \eqref{eq:surr-stoch} from Section \ref{sec:stochastic}.
Notice that $\Tilde{\ell}_t$ consists of two main parts: an inner product term that is linear in $v_{\theta} =g(\theta)$, and a squared error term $\|v_{\theta}- v_{\theta_t}\|^2$.
Our two surrogate methods use  different estimates for the linear and error term.

The first method, Surr-GD, uses the same batch of data to approximate both the inner product part and the squared error term.
This fixed batch approach was proposed in \citet{lavington2023target} for scalar minimization, and introduces bias due to approximating the squared error term across all states with states seen in the batch.
With this approximation we perform either a fixed number of steps (inner = $\#$ steps) or an adaptive number of steps depending on a loss ratio (Thresholded Surr-GD).
Taking one step recovers a batch version of TD(0), for more details see Algorithms \ref{alg:surrogate_inner},  \ref{alg:thresholded_surr}.

The second method, surrogate with double sampling, uses a fixed batch to approximate the inner product term but resamples a new batch to approximate the squared error part at each gradient step.
By resampling a new batch for the squared error we have removed the bias from the first method at the cost of increased variance in the gradient.
For more details see Algorithm \ref{alg:double_sampling}.

All methods are evaluated with value prediction error, the squared error for a test sample of 10,000 states visited by a fixed reference policy.
The values in these test states are approximated by 10 Monte-Carlo roll-outs.
As shown in Figure~\ref{fig:nonlinear}, surrogate methods significantly outperformed TD(0) in both data efficiency (outer loop iterations) and wall-clock time.
Interestingly, the surrogate methods converge to a lower prediction error than TD(0), as indicated by the red dashed line. 
% being lower than TD(0)'s error curve.

Although the inner loop in the surrogate methods requires more gradient computations, the computational overhead is marginal when compared to environment interactions, a common bottleneck in RL tasks~\citep{vaswani2021general}.
We observe that increasing the inner loop size improves data efficiency at the trade-off of a potentially slower convergence in wall-clock time. 
An inner loop count, such as $50$, seems to balance  both data efficiency and wall-clock time. 


% We compare three different surrogate methods with TD(0). 
% To demonstrate the practical relevance of our algorithm, we conducted experiments on RL tasks with nonlinear function approximation in Mujoco environments \citep{todorov2012mujoco}. These experiments compared TD(0), the standard TD learning method, with several surrogate methods designed to minimize value prediction error. The surrogate variants included one with double sampling in the inner loop to mitigate bias \citep{lavington2023target}, and another with an adaptive threshold on alpha and inner loop count.

% Each algorithm was trained using a two-layer neural network with identical initialization. To account for randomness, experiments were repeated five times with different random seeds. Value estimation was performed in an on-policy setting, with a fixed reference policy, and comparisons were made using a test set collected from this policy.

% We compared the performance of the algorithms across two dimensions; outer loop iterations, which indicates data efficiency, and wallclock time, which reflects the practical efficiency.

% As shown in Figure~\ref{fig:nonlinear}, surrogate methods significantly outperformed TD(0) in both dimensions.
% Although the inner loop in the surrogate methods requires more gradient computations,  the computational overhead is marginal comparing to the time required for environment interactions, a common bottleneck in RL tasks\citep{vaswani2021general}.

% Comparing surrogate methods with different inner loop steps, we observe that while more inner loop steps improved data efficiency, they led to slower convergence in wallclock time. A moderate inner loop count, such as $m=50$, balanced both criteria, showing that gains from additional steps become marginal as costs increase. 
% Additionally, as indicated by the red dashed line presenting the minimum value prediction error, the surrogate methods also achieved lower value prediction error on unseen test data, suggesting better generalization compared to TD(0).
\section{Conclusion}

In this work we proposed a principled and scalable surrogate loss approach to solving variational inequality problems with hidden monotone structure.
We have presented a novel $\alpha$-descent condition that is optimizer agnostic and quantifies sufficient progress on the surrogate for convergence.  
We have proved linear convergence in both the deterministic and stochastic settings.  We have also shown that surrogate losses in VI problems are strictly more difficult to analyze than in scalar minimization. 
Furthermore, we have demonstrated the generality of our approach by showing how existing  methods can be viewed as special cases of our general framework.
Empirically, we have demonstrated the effectiveness of using surrogate losses in both min-max optimization and minimizing projected Bellman error. Finally, by using surrogate losses, we have proposed novel variants of TD(0) that are more data and run-time efficient in the deep reinforcement learning setting. 


\subsubsection*{Acknowledgments}
We thank Nicolas Le Roux, David Kanaa, Iosif Sakos, Andrew Patterson, Khurram Javed, and Arushi Jain, for helpful discussions
and feedback. This work was supported by Borealis AI through the Borealis AI
Global Fellowship Award.

\bibliography{iclr}
\bibliographystyle{iclr2025_conference}
\clearpage
\appendix

\section{Proofs}
\begin{remark}\label{remark:inf}
If $g$ is continuous, and $\{g(\theta): \theta  \in\reals^d\}$ is convex with $\mathcal{Z}$ as its closure, then the least-squares surrogate loss $\ell_t(\theta)=\frac{1}{2}\|g(\theta)-z_{t}+\eta F(z_t)\|^2$ admits a unique point $z_t^\ast \in \mathcal{Z}$ such that
\[
    \ell_t^\ast =  \frac{1}{2}\|z_t^\ast-z_{t}+\eta F(z_t)\|^2,
\]
and for any $\theta$
\[
\frac{1}{2}\|g(\theta) -z_t^\ast\|^2 \leq \ell_t(\theta)-\ell_t^\ast,
\]
where $\ell_t^\ast = \inf_{\theta \in \reals^d} \ell_t(\theta)$.
\end{remark}
\begin{proof}
    Let $f(z) = \frac{1}{2}\|z-z_t+\eta F(z_t)\|^2$ be the surrogate loss with respect to the predictions $z = g(\theta)$.
    We have that $f(z) = \ell_t(\theta)$ for all $\theta \in \reals^d$.
    Now consider the set $\mathcal{Z}=\closure{\{g(\theta): \theta \in \reals^d\}}$, since it is closed and convex, we have that $f$ has a unique minimum $z_t^\ast$ because it is 1-strongly convex. 
    Furthermore, we have
    \[
        \frac{1}{2}\|z-z_t^\ast\|^2 \leq f(z) -f(z_t^\ast).
    \]
    Now since $z_t^\ast \in \mathcal{Z}$ and $\mathcal{Z}$ is the closure of $\{g(\theta): \theta \in  \reals^d\}$, there exists a sequence of parameters $\{\theta_t\}_{t\in \nnumbers}$ such that
    $\{z_t = g(\theta_t)\}_{t\in\nnumbers} \to z_t^\ast$.
    Therefore we have that,
    \begin{align*}
        f(z_t^\ast) = \lim_{t \to \infty} f(z_t) = \lim_{t \to \infty} \ell_t(\theta_t) \geq \ell_t^\ast \geq f(z_t^\ast).
    \end{align*}
    Where we have used the continuity of $f$ and $g$. The last inequality follows because $\{g(\theta): \theta \in  \reals^d\} \subseteq \mathcal{Z}$.
    Therefore $\ell_t^\ast = f(z_t^\ast)$ and the result follows.
\end{proof}

\begin{lemma}\label{thm:bias}
    If $F$ is monotone and $z_\ast$ is in the relative interior of the constraint $\mathcal{Z}$, then 
    \[
        \ell_t(\theta_t) - \ell_t^\ast \leq \frac{\eta^2 }{2}\|F(z_t)-F(z_\ast)\|^2.
    \]
    Furthermore, under the $\alpha$-descent condition (Definition \ref{eq:alpha}) we have
    \[
    \|z_{t+1}-z_t^\ast\| \leq \alpha \eta \|F(z_t)-F(z_\ast)\|.
    \]
\end{lemma}
\begin{proof}
    If $z_\ast$ is a solution then we have
    \[
        \inner{F(z_\ast)}{z-z_\ast} \geq 0\,,\; \forall z \in \mathcal{Z}.
    \]
    If $z_\ast$ is in the relative interior then for any $z \in \mathcal{Z}$ there exists a $\lambda > 1$ such that $z' = (1-\lambda)z +\lambda z_\ast  \in \mathcal{Z}$~\citep{rockafellar1997convex}[Theorem 6.4].
    Therefore by optimality we have
    \begin{align*}      
       0 &\leq \inner{F(z_\ast)}{z'-z_\ast} 
       =\inner{F(z_\ast)}{ (1-\lambda)z +\lambda z_\ast-z_\ast} 
        =(1-\lambda)\inner{F(z_\ast)}{z-z_\ast} \leq 0.   
    \end{align*}
    Where the last inequality follows from $\lambda > 1$. Altogether we have 
    \[
        \inner{F(z_\ast)}{z-z_\ast} = 0 \,,\;\forall z \in \mathcal{Z}.
    \]
    Moreover, as a consequence we have that for any two points $z, z' \in \mathcal{Z}$
    \[
        \inner{F(z_\ast)}{z-z'} = \inner{F(z_\ast)}{z-z_\ast} + \inner{F(z_\ast)}{z_\ast-z'} = 0.
    \]
    Letting $z_t^\ast$ be the exact projected gradient step and $z_t =g(\theta_t)$ the current iterate  we have
    \begin{align*}
        \ell_t(\theta_t) - \ell_t^\ast &= \frac{1}{2}\|\eta F(z_t)\|^2 - \left(\frac{1}{2}\|z_t^\ast -z_t + \eta F(z_t)\|^2 \right)\\
        &=\eta \inner{F(z_t)}{z_t-z_t^\ast}  -\frac{1}{2}\|z_t^\ast-z_t\|^2\\
        &= \eta \inner{F(z_t) - F(z_\ast)}{z_t-z_t^\ast} + \eta \inner{F(z_\ast)}{z_t-z_t^\ast}  -\frac{1}{2}\|z_t^\ast-z_t\|^2\\
        &= \eta \inner{F(z_t) - F(z_\ast)}{z_t-z_t^\ast} -\frac{1}{2}\|z_t^\ast-z_t\|^2\\
        &\leq \frac{\eta^2 }{2}\|F(z_t)-F(z_\ast)\|^2.
    \end{align*}
    Where the last two inequalities follow by: $z_\ast$ being within the relative interior, and the inequality $\inner{u}{v} \leq \frac{\rho}{2}\|u\|^2 +\frac{1}{2\rho}\|v\|^2$ for any $\rho > 0$.

    By Remark \ref{remark:inf} and the $\alpha$-decent condition,
    \[
        \frac{1}{2}\|z_{t+1}-z_t^\ast\|^2 \leq \ell_t(\theta_{t+1})-\ell_t^\ast \leq \alpha^2 (\ell_t(\theta_t) -\ell_t^\ast) \leq \alpha^2\frac{\eta^2}{2}\|F(z_t)-F(z_\ast)\|^2.
    \] 
\end{proof}


\begin{proposition}\label{thm:alpha-scalar}
    Let $f: \reals^n \to \reals$ be $L$-smooth.
    For some $g: \reals^d \to \reals^n$ define the surrogate loss $\ell_t(\theta) = \frac{1}{2}\|g(\theta)-z_t+\frac{1}{L}\nabla f(z_t)\|^2$,
    where $z_t = g(\theta_t)$.
    Then the $\alpha$-descent condition (Definition \ref{eq:alpha}) guarantees
    \[
    f(z_{t+1}) \leq f(z_t) -L(1-\alpha^2)(\ell_t(\theta_t)- \ell_t^\ast).
    \]
    If $f$ is bounded below then $\{z_t - z_{t}^\ast\}_{t\in \nnumbers}\to 0$, where $z_t^\ast = \Pi(z_t - \frac{1}{L}\nabla f(z_t))$.
\end{proposition}
\begin{proof}
    Let $\hat{z_t} = z_t- \frac{1}{L}\nabla f(z_t)$.
    \begin{align*}
        f(z_{t+1}) - f(z_t) &\leq \inner{\nabla f(z_t)}{z_{t+1}-z_t} + \frac{L}{2}\|z_{t+1}-z_t\|^2\\
        &= L\inner{z_t-\hat{z}_t}{z_{t+1}-z_t}+ \frac{L}{2}\|z_{t+1}-z_t\|^2\\
        &=L\left(-\|\hat{z}_t - z_t\|^2 +\inner{z_t-\hat{z}_t}{z_{t+1}-\hat{z}_t} + \frac{1}{2}\|z_{t+1}-z_t\|^2\right)\\
        &=L\left(\frac{1}{2}\|\hat{z}_t-z_{t+1}\|^2-\frac{1}{2}\|\hat{z}_t-z_t\|^2 \right)\\
        &=L\left(\ell_t(\theta_{t+1})-\ell_{t}(\theta_t) \right).
    \end{align*}
    The second to last equality follows from  expanding $\frac{1}{2}\|z_{t+1} -z_t\|^2 = \frac{1}{2}\|z_{t+1}-\hat{z}_t +\hat{z}_t- z_t\|^2$.
    Using the $\alpha$-descent condition we have
    \[
    \ell_t(\theta_{t+1}) -\ell_{t}(\theta_t) \leq (\alpha^2-1)( \ell_t(\theta_t) -\ell_t^\ast),
    \]
    yielding the first result.

    If $f$ is bounded below by some constant $c$ then we have:
    \begin{align*}
        \sum_{t=1}^T L(1-\alpha^2)(\ell_t(\theta_t)- \ell_t^\ast) \leq f(z_1)-f(z_T) \leq f(z_1)-c.
    \end{align*}
    Therefore the series $\sum_{t=1}^\infty\ell_t(\theta_t)- \ell_t^\ast$ converges with $\ell_t(\theta_t)- \ell_t^\ast \to 0$.
    By Remark \ref{remark:inf} we have  $\frac{1}{2}\|z_t - z_t^\ast\|^2 \leq \ell_t(\theta_t)- \ell_t^\ast$ implying the result $z_t - z_t^\ast \to 0$.
\end{proof}


{
\renewcommand{\thetheorem}{\ref{thm:det}}
\begin{theorem}
    Let Assumption \ref{assumption:vi} hold and let $\{z_t = g(\theta_t)\}_{t\in \nnumbers}$ be  the iterates produced by Algorithm~\ref{alg:surr}. If $\alpha$ and $\eta$ are picked  such that
$
\rho:= 1-2\eta(\mu-\alpha L) + (1+\alpha^2)\eta^2 L^2 <1 
$ then, $z_t$ converge linearly to the solution $z_\ast$ at the following linear rate:
\begin{equation}
    \|z_{t+1} -z_*\|^2 \leq \rho^t \|z_1-z_*\|^2.
\end{equation}Particularly, if $\alpha< \frac{\mu}{L}$ and $\eta < \frac{2 (\mu-\alpha L)}{(1+\alpha^2)L^2}$ then $\rho <1$ and if
$\alpha \leq \frac{\mu}{2L}$ and $\eta = \frac{2\mu}{5L^2}$ then $\rho \leq 1 - \frac{\mu^2}{5L^2}$.
\end{theorem}
}
\begin{proof}
First note that by definition of Algorithm \ref{alg:surr}, the iterates ${z_t = g(\theta_t)}_{t\in\nnumbers}$ satisfy the $\alpha$-descent property (Definition \ref{eq:alpha}), therefore Lemma \ref{thm:bias} holds.
Recall that $z_t^\ast$, the exact projection update, is a contraction if $\eta < \frac{2\mu}{L}$ since 
\begin{equation}
    \label{eq:proj:contract}
    \|z_t^\ast -z_\ast\|^2 \leq \kappa^2\|z_t-z_\ast\|^2
\end{equation}
where $\kappa^2 = 1-2\eta \mu + \eta^2 L^2$\cite[Theorem 12.1.2]{facchinei2003finite}.
For the remainder, assume that $\eta < \frac{2\mu}{L}$ so that $\kappa \in [0, 1)$.
Denoting $F_t = F(z_t)$ and $F_\ast = F(z_\ast)$, we have
    \begin{align*}
        \|z_{t+1}-z_\ast\|^2 &= \|z_{t}^\ast -z_\ast + z_{t+1} - z_{t}^\ast \|^2\\
        &= \|z_t^\ast-z_\ast\|^2 + 2\inner{z_{t+1}-z_t^\ast}{z_t^\ast-z_\ast} + \|z_{t+1} - z_{t}^\ast \|^2\\
        &\leq \|z_t^\ast-z_\ast\|^2 + 2\|z_{t+1}-z_t^\ast\|\|z_t^\ast-z_\ast\| + \|z_{t+1} - z_{t}^\ast \|^2 \qquad\qquad \text{(Cauchy–Schwarz)}\\
        &\leq \|z_t^\ast-z_\ast\|^2 + 2\alpha \eta \kappa  \|F_t- F_\ast\|\|z_t-z_\ast\| + \alpha^2 \eta^2\|F_t - F_\ast\|^2 \qquad\qquad \text{(Lemma~\ref{thm:bias})}\\
        &\leq  \|z_t^\ast-z_\ast\|^2 + 2\alpha \eta L\|z_t-z_\ast\|^2 + \alpha^2 \eta^2L^2\|z_t-z_\ast\|^2\;\; \text{(Smoothness of $F$ and $\kappa<1$)}\\
        &\leq \kappa^2 \|z_t-z_\ast\|^2 + 2\alpha \eta L\|z_t-z_\ast\|^2 + \alpha^2 \eta^2L^2\|z_t-z_\ast\|^2\qquad \qquad \qquad \qquad \;\, \text{(Eq.~\ref{eq:proj:contract})}\\
        &= \|z_t-z_\ast\|^2\left(1-2\eta \mu +2\alpha \eta L + (1+\alpha^2)\eta^2 L^2\right).
    \end{align*}
If $\alpha < \frac{\mu}{L}$ then
\begin{align*}
  \|z_{t+1}-z_\ast\|^2 \leq \|z_{t}-z_\ast\|^2\left( 1 -2\eta \underbrace{(\mu -\alpha L)}_{>0} + (1+\alpha^2)\eta^2 L^2\right).
\end{align*}
Taking $\eta < \frac{2(\mu-\alpha L)}{(1+\alpha^2)L^2}$ would guarantee a contraction.

If $\alpha \leq \frac{\mu}{2L}$ and taking $\eta = \frac{2\mu}{5L^2}$ we have: 
\begin{align*}
    1-2\eta \mu +2\alpha \eta L + (1+\alpha^2)\eta^2 L^2 &\overset{\alpha \leq \nicefrac{\mu}{2L} \leq \nicefrac{1}{2}}{\leq} 1-\eta \mu +\left(1+\frac{1}{4}\right)\eta^2 L^2 \\
    &\overset{\eta = \frac{2\mu}{L^2}}{=}1-\frac{2\mu^2}{5L^2} +\frac{\mu^2}{5L^2} = 1- \frac{\mu^2}{5L^2}.
\end{align*}


\end{proof}

\begin{proposition}\label{thm:contraction-error-limit}
Take $\theta_{t+1}$ to be an approximate minima of $\ell_t(\theta)$. 
Suppose $T = \Pi \circ (\id - \eta F)$ is a contraction, if $\ell_t(\theta_{t+1})- \ell_t^\ast \to 0$ then the induced sequence $\{z_{t} = g(\theta_t)\}_{t \in \nnumbers}$ converges, $z_t \to z_{\ast}$ where $z_{\ast}$ is the unique solution to $\vi(\mathcal{Z}, F)$.
\end{proposition}

\begin{proof}
Let $\epsilon_t = z_{t+1}-z_t^\ast$  be the approximation error between $z_{t+1}$ and $z_t^\ast$ the minimum  of the surrogate $\ell_t$ (exact projected gradient step). 
    \begin{align*}
        \frac{1}{2}\| z_{t+1}-z_\ast\|^2 &= \frac{1}{2}\| z_{t}^\ast-z_\ast + \epsilon_t\|^2\\
        &= \frac{1}{2}\| z_{t}^\ast-z_\ast\|^2 + \langle z_{t}^\ast-z_\ast, \epsilon_t\rangle + \frac{1}{2}\|\epsilon_t\|^2\\
        \overset{\rho > 0}&{\leq} \frac{1}{2}\| z_{t}^\ast-z_\ast\|^2+\frac{\rho}{2}\| z_{t}^\ast-z_\ast\|^2 +\frac{1}{2\rho}\|\epsilon_t\|^2 + \frac{1}{2}\|\epsilon_t\|^2\\
        \overset{\kappa \in [0,1)}&{\leq} \frac{\kappa(1 + \rho)}{2} \| z_{t}-z_\ast\|^2 + (1+\frac{1}{\rho})(\ell_t(\theta_{t+1})-\ell_t^\ast).
        \end{align*}
Where we use the fact that $\frac{1}{2}\|\epsilon_t\|^2 = \frac{1}{2}\|z_{t+1}-z_t^\ast\|^2 \leq \ell_t(\theta_{t+1})-\ell_t^\ast$ from Remark \ref{remark:inf}.
Take any $\rho$ such that $\kappa(1 + \rho) <1$ then apply Lemma 3.9 in \cite{FRANCI2022161}.
\end{proof}
{
\renewcommand{\thetheorem}{\ref{prop:counter}}
\begin{proposition}
There exists an $L$-smooth and $\mu$-strongly monotone $F$, and a sequence of iterates $\{z_t\}_{t \in \nnumbers}$ verifying the alpha descent condition with  $\alpha < 1$  such that  $z_t$ diverges for any $\eta$.
\end{proposition}
}
\begin{proof}
    

Let us consider the simple min max example with loss
\[
    \min_x \max_y \frac{1}{2}x^2 + xy -\frac{1}{2}y^2.
\]
This problem can be written as VIP where $z = (x,y)$ and the operator $F$ is linear given by the matrix 
\[
F = \begin{bmatrix}
    1 & 1 \\
    -1 & 1
\end{bmatrix}.
\]
It is well-known that $F$ is both smooth and strongly monotone with $z_{t+1} = z_t -\eta F(z_t)$ being a contraction for small enough $\eta$~\citep{facchinei2003finite}.
Now if we consider a biased direction given by a matrix $P$, that is $z_{t+1} = z_t -P z_t$, then using the fact that $\ell_t^\ast =0$ the $\alpha$-descent on the surrogate corresponds to the following bound
\begin{align*}
    \|z_{t+1}-z_t + \eta F z_t\| = \|(\eta F -P) z_t\| \leq \alpha \|\eta F z_t\|.
\end{align*}
Despite being a contraction when we follow the true gradient $F$, the above min max loss causes rotations in the dynamics that are inherent to the adversarial nature of the problem. These rotations are carefully controlled by the stepsize and strong convexity/concavity of the loss.
Our counterexample simply adds a bit of rotation that ensures $\alpha <1 $ but yet is detrimental to the convergence.
Mathematically, we take $P = (\id -\alpha Q) \eta F$ where $Q$ is the rotation matrix
\[
Q = 
\begin{bmatrix}
    \frac{1}{\sqrt{2}}  & -\frac{1}{\sqrt{2}} \\
    \frac{1}{\sqrt{2}}  & \frac{1}{\sqrt{2}}. 
\end{bmatrix}
\]
With this construction we are guaranteed that $\alpha < 1$ since
\begin{align*}
    \|(\eta F -P) z_t\| = \|\alpha Q \eta F z_t\| =  \alpha\eta \|F z_t\|,
\end{align*}
where the last equality is due the the fact that $Q$ is an orthogonal matrix and therefore does not change the Euclidean norm of a vector.

Now taking $\alpha =\nicefrac{1}{\sqrt{2}}$ gives 
\begin{align*}
    P &= \left(\id - \alpha\begin{bmatrix}
    \frac{1}{\sqrt{2}}  & -\frac{1}{\sqrt{2}} \\
    \frac{1}{\sqrt{2}}  & \frac{1}{\sqrt{2}} 
\end{bmatrix}\right)
\begin{bmatrix}
    \eta & \eta \\
    -\eta & \eta
\end{bmatrix}.\\
&= \begin{bmatrix}
    \frac{1}{2}  & \frac{1}{2} \\
    -\frac{1}{2}  & \frac{1}{2} 
\end{bmatrix}
\begin{bmatrix}
    \eta & \eta \\
    -\eta & \eta
\end{bmatrix}=\begin{bmatrix}
    0 & \eta \\
    -\eta & 0.
\end{bmatrix}
\end{align*}
We have that $z_{t+1} = z_t - Pz_t = \begin{bmatrix}
    1 & -\eta \\
    \eta & 1
\end{bmatrix}z_t$ has an $\alpha = \nicefrac{1}{\sqrt{2}}$ but yet diverges for any $\eta > 0$ since the Eigenvalues of the linear system are $\lambda = 1 \pm i\eta$ therefore the spectral radius is strictly greater than one. 
Note that these dynamics are equivalent to gradient descent ascent on the bilinear loss $f(x,y) = xy$, which is known to diverge for any stepsize.

\end{proof}

{
\renewcommand{\thetheorem}{\ref{thm:stochastic}}
\begin{theorem}
Let $\mathcal{Z} = \reals^n$, and Assumption~(\ref{eq:exp-co}) hold.  
If $F_{\xi}(x)$  is $L$-smooth and $\eta \leq \frac{1}{2(1+c)L}$ where  $c \geq  2(1+\alpha^2)$  then any trajectory $\{\theta_t\}_{t\in \nnumbers}$ 
satisfying the $\alpha$-expected descent condition guarantees:
\[
\expect{\tfrac{1}{2}\|z_{t+1}-z_\ast\|^2} \leq \expect{\tfrac{1}{2}\|z_{t}-z_\ast\|^2}\left(1-\eta\mu + \alpha^2\right) +\eta^2(1+c)\sigma^2. 
\]
\end{theorem}
}
\begin{proof}
For convenience let $F_{\xi_t}(z_t) = \hat{F}_t$.
Let $z_t^\ast = z_t -\eta \hat{F}_t$ and $p_t = z_t-z_{t+1}$, $\expect[\theta_{t+1}]{\cdot}$ denote the expectation over $z_{t+1} =g(\theta_{t+1})$ given $z_t$, and $\xi_t$.
The expected $\alpha$-descent condition in the unconstrained case is equivalent to
\[
    \expect[\theta_{t+1}]{\|z_{t+1}-z_t + \eta \hat{F}_t\|^2} \leq \alpha^2 \eta^2 \|\hat{F}_t\|^2.
\]
By concavity of the square root and the  expected $\alpha$-descent condition we have $\expect{\sqrt{X}} \leq \sqrt{\expect{X}}$ and
\begin{align}
    \expect[\theta_{t+1}]{\|z_{t+1}-z_t +\eta \hat{F}_t\|} &= \expect[\theta_{t+1}]{\sqrt{\|z_{t+1}-z_t +\eta \hat{F}_t\|^2}} \\
    &\leq \sqrt{\expect{\|z_{t+1}-z_t +\eta \hat{F}_t\|^2}} \\
    &\leq \alpha \eta \|\hat{F}_t\|.\label{stoch_ineq_1}
\end{align}
We also have the following bound
\begin{align}
    \expect[\theta_{t+1}]{\|z_{t+1}-z_t\|^2} &= \expect[\theta_{t+1}]{\|z_{t+1}-z_t + \eta \hat{F}_t -  \eta \hat{F}_t\|^2}  \\ &\leq \expect[\theta_{t+1}]{2\|z_{t+1}-z_t + \eta \hat{F}_t\|^2 + 2\|\eta \hat{F}_t\|^2}\\
    &\leq 2(1+\alpha^2)\eta^2\|\hat{F}_t\|^2.\label{stoch_ineq_2}
\end{align}

Next we expand and bound the distance $\frac{1}{2}\|z_{t+1}-z_\ast\|^2$.

    \begin{align*}
   \frac{1}{2}\|z_{t+1}-z_\ast\|^2 &=\frac{1}{2}\| z_{t}-p_t-z_\ast\|^2 \\
    &= \frac{1}{2}\|z_{t}-z_\ast\|^2 -\inner{p_t}{z_t-z_\ast} + \frac{1}{2}\|p_t\|^2\\
    &=\frac{1}{2}\|z_{t}-z_\ast\|^2 -\eta\inner{\hat{F}_t}{z_t-z_\ast} +\inner{\eta \hat{F}_t-p_t}{z_t-z_\ast} + \frac{1}{2}\|p_t\|^2\\
    &\leq \frac{1}{2}\|z_{t}-z_\ast\|^2 -\eta\inner{\hat{F}_t}{z_t-z_\ast} +\|\eta \hat{F}_t-p_t\|\|z_t-z_\ast\| + \frac{1}{2}\|p_t\|^2\\
    &= \frac{1}{2}\|z_{t}-z_\ast\|^2 -\eta\inner{\hat{F}_t}{z_t-z_\ast} +\|z_{t+1}-z_t+ \eta \hat{F}_t\|\|z_t-z_\ast\| + \frac{1}{2}\|z_{t+1}-z_t\|^2
    \end{align*}
Taking an expectation over $z_{t+1}$ given $z_t$, and $\xi_t$ we have 
    \begin{align*}
    \expect[\theta_{t+1}]{\frac{1}{2}\|z_{t+1}-z_\ast\|^2 } \overset{(\ref{stoch_ineq_1}, \ref{stoch_ineq_2})}&{\leq} \frac{1}{2}\|z_{t}-z_\ast\|^2 -\eta\inner{\hat{F}_t}{z_t-z_\ast} +\eta \alpha \|\hat{F}_t\|\|z_t-z_\ast\|+ \frac{2(1+\alpha^2) \eta^2}{2}\|\hat{F}_t\|^2\\
    &\leq \frac{1}{2}\|z_{t}-z_\ast\|^2 -\eta\inner{\hat{F}_t}{z_t-z_\ast} +\eta \alpha \|\hat{F}_t\|\|z_t-z_\ast\|+ \frac{c \eta^2}{2}\|\hat{F}_t\|^2\\
    &\leq \frac{1}{2}\|z_{t}-z_\ast\|^2 -\eta\inner{\hat{F}_t}{z_t-z_\ast} +\frac{\alpha^2}{2} \|z_t-z_\ast\|^2+ \frac{\eta^2(1+c)}{2}\|\hat{F}_t\|^2
\end{align*}
where in the inequality we used (\ref{stoch_ineq_1}, \ref{stoch_ineq_2}) and in the second inequality $c \geq 2(1+\alpha^2)$.
Taking an expectation given information up to time $t$  and using the expected co-coercivity assumption gives,
\begin{align*}
    \expect[t]{\frac{1}{2}\|z_{t+1}-z_\ast\|^2} &\leq \frac{1}{2}\|z_{t} -z_\ast\|^2 -\eta\inner{F_t}{z_t-z_\ast} +\frac{\alpha^2}{2}\|z_t-z_\ast\|^2+ \frac{\eta^2(1+c)}{2} \expect[t]{\|\hat{F}_t\|^2 }\\
    \overset{\eqref{eq:exp-co}}&{\leq} \frac{1}{2}\|z_{t}-z_\ast\|^2 -\eta\inner{F_t}{z_t-z_\ast} + \frac{\alpha^2}{2}\|z_t-z_\ast\|^2+ \eta^2(1+c) L\inner{F_t}{z_t-z_\ast} + \eta^2(1+c)\sigma^2 \\
    &= \frac{1}{2}\|z_{t}-z_\ast\|^2 -\eta\left(1-\eta(1+c)L\right)\inner{F_t}{z_t-z_\ast} + \frac{\alpha^2}{2}\|z_t-z_\ast\|^2+ \eta^2(1+c)\sigma^2 \\ 
    \overset{\eta (1+c)L< 1}&{\leq} \frac{1}{2}\|z_{t}-z_\ast\|^2 -\eta\mu\left(1-\eta(1+c)L\right)\|z_t-z_\ast\|^2 + \frac{\alpha^2}{2}\|z_t-z_\ast\|^2+ \eta^2(1+c)\sigma^2 
\end{align*}
If we take $\eta < \frac{1}{2(1+c)L}$ then we have the following recursion
\[
  \expect[t]{\frac{1}{2}\|z_{t+1}-z_\ast\|^2} \leq \frac{1}{2}\|z_{t}-z_\ast\|^2\left(1-\eta \mu + \alpha^2\right) + \eta^2(1+c)\sigma^2.
\]
Taking an expectation over $z_t$ gives the result.
\end{proof}


\section{Non-linear least-squares}
% \begin{proposition}
%     If $g: \reals^d \to \reals^n$ is continuously differentiable and the transpose of its Jacobian $Dg(\theta)^\top$ has globally lower-bounded singular values, then the its range $\mathcal{Z} = \{g(\theta): \theta \in \reals^d\}$ is unbounded.
% \end{proposition}
% \begin{proof}
%     We show that the following function $f(t) = \inner{g(tu)}{v}$ for any $\|u\|_2 =1$ and $\|v\|_2=1$ is unbounded. Which is sufficient since
%     \[
%     \|f(t)\| \leq \|g(tu)\|\|v\|= \|g(tu)\|.
%     \]
%     Therefore $\lim_{t\to +\infty}f(t) = +\infty \Rightarrow \lim_{t\to +\infty}\|g(tu)\| = +\infty$.
% First we lower bound the derivative of $f$. 
% we use $\|x\|_\infty = max_i|x^i|$ to denote the infinity norm, and recall that the following bounds hold $\|x\|_\infty \leq\|x\|_2 \leq \sqrt{d}\|x\|_\infty$. 
% By chain rule we have that $f'(t) = \inner{u}{Dg(tu)^\top v}$ and the following:
% \begin{align}
%     \|f'(t)\| &= \|\inner{u}{Dg(tu)^\top v}\|\\
%     &=\|\sum_{i=1}^d u^i (Dg(tu)^\top v)^i\|
%     &\geq \max_{i}
% \end{align}
% \end{proof}

\begin{lemma}\label{thm:smooth-surr}
Let $f(\theta) = \frac{1}{2}\|r(\theta)\|^2$ where $r: \reals^d \to \reals^n$, $r(\theta) = (r_i(\theta),\cdots, r_n(\theta))^\top$.
    Suppose  $r_i(\theta)$ is $\beta$-smooth and $Dr(\theta)^\top$ has singular values that are globally upper bounded by $\sigma_{\max} >0$.
    If $r$ is bounded,  then $f$ is $L$-smooth: there exists $L \geq 0$  such that $\|\nabla f(\theta) -\nabla f(\theta')\| \leq L\|\theta -\theta'\|$ for any $\theta, \theta' \in\reals^d$.
\end{lemma}

\begin{proof}
Since each $r_i$ is $\beta$-smooth we have that
    \begin{align*}
    \|Dr(\theta)^\top-Dr(\theta')^\top\|_{1,2} &= \sup\{\|(Dr(\theta)^\top-Dr(\theta')^\top)v\|_{1}: v \in \reals^n , \|v\|_2\leq 1\} \\ 
    &=\max_{i}\|\nabla g_i(\theta)-\nabla g_i(\theta')\| \leq L\|\theta-\theta'\|.    
    \end{align*}
        % 
    \begin{align*}
        \|\nabla f(\theta)-\nabla f(\theta')\| &= \|Dr(\theta)^\top r(\theta) - Dr(\theta')^\top r(\theta')\| \\
        &= \|Dr(\theta)^\top r(\theta) -Dr(\theta)^\top r(\theta')+ Dr(\theta)^\top r(\theta') - Dr(\theta')^\top r(\theta')\|\\
        &\leq \|Dr(\theta)^\top r(\theta) -Dr(\theta)^\top r(\theta')\| + \|Dr(\theta)^\top r(\theta') - Dr(\theta')^\top r(\theta')\|\\
        &\leq \sigma_{\max}\|r(\theta)-r(\theta')\| + \|Dr(\theta)^\top - Dr(\theta')^\top\|_{1,2}\|r(\theta')\|_1\\
        &\leq \sigma_{\max}^2\|\theta - \theta'\|+\beta C\|\theta-\theta'\|.
    \end{align*}
    Where the last inequality follows from $r$ being $\sigma_{\max}$ Lipschitz since $\|r(\theta) - r(\theta')\| \leq \|Dr(\theta)\|\|\theta-\theta'\|$ by the mean value inequality~\citep{hormander2007analysis} and by the fact that $\|Dr(\theta)^\top\| \leq \sigma_{\max}$.
    Since $r$ is bounded it follows that there exists $\|r(\theta)\|_1 \leq C$ for all $\theta$.
\end{proof}

\begin{definition}[\cite{polyak1964gradient,lojasiewicz1963propriete}]\label{def:pl}
A function $f: \reals^n \to \reals$ 
satisfies the PL condition if there exists $\mu > 0$
such that for all $x \in \reals^n$
\begin{align}
   \|\nabla f(x)\|^2 \geq 2 \mu (f(x) -f^\ast),
\end{align}
where $f^\ast = \inf_{x \in \reals^n}f(x)$. 
\end{definition}



{
\renewcommand{\thetheorem}{\ref{thm:pl}}
\begin{proposition}
    Assume $f$ satisfies the $\mu$-PL condition, and let $\sigma_{\min}$ be a lower bound on the singular values of $Dg(\theta)^\top$. Then, we have $f \circ g$ is $\mu \sigma^2_{\min}$-PL.
\end{proposition}
}
\begin{proof}
We have that 
\begin{align*}
 || \nabla f \circ g (\theta)||^2 &= ||Dg(\theta)^\top \nabla_{z=g(\theta)} f(z)||^2 \\
 &=\nabla f(z)^\top D(g(\theta)) D(g(\theta))^\top \nabla f(z)\\
 &\geq \sigma_{\min}^2||\nabla f(z)||^2   \\
 &\geq \sigma_{\min}^2 \mu(f(z)-\inf_{z'} f(z')) \\
    &\geq \sigma_{\min}^2 \mu(f(z)-\min_{z' \in \mathcal{Z}} f(z')) \\
 &= \sigma_{\min}^2 \mu (f(g(\theta))- \inf_{\theta'}f(g(\theta'))).
\end{align*}
\end{proof}

\begin{proposition}
    If $g: \reals^d \to \reals^n$ $g(\theta) = (g_1(\theta), \cdots, g_n(\theta))^\top$ is differentiable where $g_i$ is $\beta$-smooth and $Dg(\theta)^\top$ has globally lower and upper bounded singular values, then  $\{g(\theta): \theta\in \reals^d\}$ is unbounded.
\end{proposition}
\begin{proof}
We prove by contradiction. Suppose $\{g(\theta): \theta\in \reals^d\}$ is bounded and denote $\mathcal{Z}$ as its closure which is also bounded. Then there exists $v \in \reals^n \notin \mathcal{Z}$. 
Consider the function $f(\theta) = \frac{1}{2}\|g(\theta)-v\|^2$. 
Taking $r(\theta) = g(\theta) -v$, we have $Dr = Dg$, and $r_i$ are still $\beta$-smooth, therefore by Lemma \ref{thm:smooth-surr} $f$ is $L$-smooth. Since $f$ is continuous and $\mathcal{Z}$ is compact, its minimum is attained, and thus there exists $z_\ast \in \mathcal{Z}$ such that
\[
    \frac{1}{2}\|z_\ast-v\|^2 = \inf_{\theta \in \reals^d}f(\theta).
\]
Moreover, there exists a sequence $\{\theta_t\}_{t\in\nnumbers}$ such that $g(\theta_t) \to z_\ast$ since $\mathcal{Z}$ is closed.
By the smoothness of $f$ it also follows that 
\[
    \frac{\|\nabla f(\theta_t)\|^2}{2L} \leq f(\theta_t) -\inf_{\theta \in \reals^d}f(\theta).
\]
Since singular values of $Dg(\theta)^\top$ are lower-bounded, there exists $\sigma_{\min} > 0$ such that $\|Dg(\theta)^\top v\| \geq \sigma_{\min}\|v\|$.
Therefore we have
\begin{align*}
    0 = \lim_{t}f(\theta_t) -\inf_{\theta \in \reals^d}f(\theta) &\geq   \lim_t\frac{\|\nabla f(\theta_t)\|^2}{2L} \\ 
    &= \lim_t\frac{\|Dg(\theta_t)^\top (g(\theta_t)-v)\|^2}{2L} \geq \lim_t\frac{\sigma_{\min}^2 }{2L}\|g(\theta_t) -v\|^2,
\end{align*}
which implies that $z_\ast =\lim_t g(\theta_t) =v$ and that $v$ is indeed within $\mathcal{Z}$ a contradiction!


\end{proof}



\section{Min max experiments}


\begin{figure}
    \centering
    \includegraphics[width=0.5\linewidth]{figures/time_matching_pennies.pdf}
    \caption{Time in seconds for performing 10,000 iterations of each method. The number in parenthesis correspond to number of inner steps taken. As a special case we have PHGD and GDA are equivalent to GN(1) and Surr-GD(1) respectively.}
    \label{fig:pennies-time}
\end{figure}

\begin{figure}
    \centering
    \includegraphics[width=0.8\linewidth]{figures/pennies_rand_dist.pdf}
    \caption{Average convergence for various surrogate methods from Section \ref{sec:non-linear} in the hidden matching pennies game. 
    The average was computed over 1,000 trajectories each with randomly sampled parameters $\alpha_j^i$ and $\theta_0$.
    }
    \label{fig:pennies-rand}
\end{figure}

\subsection{Hidden Matching Pennies}\label{sec:appendix-pennies}

In the hidden matching pennies game each player $i$ has the parameterization $h^i(\theta) = \sigmoid(\alpha_2^i\celu(\alpha^i_1 \theta))$.
The parameters for each player are:
\begin{align*}
    &\alpha_1^1 = 0.5 , \text{ and } \alpha_2^1 = 1 \text{ for player } 1,\\
    &\alpha_1^2 = 0.7 , \text{ and } \alpha_2^2 = 1 \text{ for player } 2.
\end{align*}
The parameters $\alpha_1^i, \alpha_2^i$, were chosen to approximately replicate the trajectory of PHGD presented in \citet[Figure 4]{sakos2024exploiting}.
The initial parameters $\theta_{1}$ were selected to be the same as those selected by \citet{sakos2024exploiting}, $\theta_1 = (\theta_1^1, \theta_1^2) = (1.25, 2.25)$.

In Figure \ref{fig:pennies-time} we report the runtime for different surrogate algorithms and number of inner steps.
In Figure \ref{fig:pennies-rand}, the average distance across 100 random initializations is observered. 
$\theta_0$ was randomly sampled from two independent Gaussians with  standard deviation of $4$. Similarly, each $\alpha_j^i$ was randomly sampled from a uniform random variable with support $[-1, 1]$.


\begin{figure}
    \centering
    \includegraphics[width=0.75\linewidth]{figures/rps_rand_dist.pdf}
    \caption{Average convergence for various surrogate methods from Section \ref{sec:non-linear} in the hidden rock-paper-scissors game. 
    The average was computed over 1,000 trajectories each with randomly sampled parameters $A_j^i$ and $\theta_0$.}
    \label{fig:rps-rand}
\end{figure}

\subsection{Hidden Rock-Paper-Scissors}\label{sec:appendix-rps}
Each player i's mixed strategy in the hidden rock-paper-scissors game is parameterized by the function $h^i(\theta) = \softmax(A_2^i \celu(A_1^i \theta^i))$, with  $\theta^i \in \reals^5$.
The min-max objective  is $f(z^1,z^2) = \inner{z^1}{A z^2} + \frac{0.2}{2}\|z-z_\ast\|^2$, where $z = (z^1,z^2)$ and $z_\ast$ is the equilibrium of rock-paper-scissors (sample rock, paper, and scissors with equal probability)~\citep{sakos2024exploiting}.
In Figure \ref{fig:rps-rand} we observe the average convergence across 100 random initializations of $\theta_1 \in \reals^5$ and matrixes $A_j^i$ for various surrogate based methods.
$\theta_1$ was sampled from a 5-dimensional isotropic Gaussian while each element of $A_j^i$ was sampled from a uniform random variable with support $[-1,1]$.

\section{Policy Evaluation and Projected Bellman Error}\label{sec:appendix-PBE}

In this section we provide some background on minimizing projected Bellman error (PBE) and the stochastic methods proposed by \citet{bertsekas2009projected}.
For a given MDP with $n$ states $(s^1, \cdots, s^n)$, and policy $\pi$, we denote the expected reward $r_\pi \in \reals^n$, discount factor $\gamma \in (0,1)$, and state to state probability transition matrix $P_\pi \in \reals^{n \times n}$.
The associated linear Bellman operator is $T_{\pi}(z) = r_{\pi} +\gamma P_{\pi}z$ and is contraction~\cite{bertsekas2012dynamic}. The value vector $v_{\pi} \in \reals^n$ is the unique fixed point of $T_{\pi}$, with i\textsuperscript{th} component $v^i$ equals to the discounted return starting from state $s^i$, $\expect[\tau = (s_1, s_2, \cdots)]{\sum_{t=1}^\infty \gamma ^t r_{\pi}(s_t)|s_1 = s^i}$, where we use $r_{\pi}(s)$ to mean $r_\pi^i$ such that $s^i = s$, and the expectation is over possible trajectories $\tau = (s_1, s_2, \cdots)$.

Under function approximation it may not be possible to find a fixed point of $T_{\pi}$ and find the true value function.
Instead we can consider the projected fixed point,
\begin{align}\label{eq:proj-fixed-point}
    z_{\ast} = \Pi_{\Xi}(T_{\pi}(z_\ast)).
\end{align}

Where the projection is on to the set of realizable value functions $\mathcal{Z} = \{v_\theta: \theta \in \reals^d\}$. 
We use the notation $v_{\theta}$ in place of $g(\theta)$ to be consistent with the RL literature, the prediction for state $s^i$ is $v_{\theta}(s^i) = g(\theta)^i$. 
$\Xi$ is the diagonal matrix with entries corresponding to the stationary distribution $\xi = (\xi^1,\cdots, \xi^n)$ of $P_{\pi}$.

As pointed out by \citet{bertsekas2009projected}, finding the fixed point \eqref{eq:proj-fixed-point} is equivalent to solving a VI of the form \eqref{eq:vi}
with constraint $\mathcal{Z}$ and operator \begin{align}
    F(z) = \Xi(z-T_\pi(z)).\label{eq:pbe-vi}
\end{align}
Additionally, \citet{bertsekas2009projected} showed that $F$ is both smooth and strongly monotone and therefore satisfies our hidden structure Assumption \ref{assumption:vi}.

This fixed point can then be equivalently mapped to the minimum of the projected Bellman error with respect to the parameters, i.e. a TD fixed point. 
The fixed point $z_\ast =v_{\theta_\ast}$ corresponds to finding parameters $\theta_\ast$ such that
\begin{align}\label{eq:pbe-loss-fixed-point}
    \theta_\ast \in \underset{\theta}{\arg\min} \left[ BE(\theta, \theta_\ast) =  \frac{1}{2} \sum_{i=1}^n \xi^i(v_{\theta}(s^i) - (r_\pi(s^i)- \gamma\expect[s'\sim P_\pi]{v_{\theta_\ast}(s')|s^i})^2\right].
\end{align}
Note this is not a scalar minimization problem since the loss is not directly minimized through the next state prediction, doing so would result in minimizing the Bellman error, $\min_{\theta} BE(\theta, \theta)$ and not PBE. 
The fixed points of these objectives need not coincide even with linear function approximation~\citep{sutton2018reinforcement}.

Using the function $BE$ and the stationary condition \eqref{eq:pbe-loss-fixed-point} we can define the following natural gap function
\begin{align}\label{eq:bellman-gap}
    BE_{\text{gap}}(\theta) = BE(\theta, \theta) - \inf_{\theta'} BE(\theta', 
    \theta).
\end{align}
Notice that $BE_{\text{gap}}(\theta) \geq 0 $ for all $\theta$ and is zero if and only if $\theta$ is a fixed point of \eqref{eq:pbe-loss-fixed-point}.



Since $\Pi_{\Xi} \circ T_{\pi}$ is known to be a contraction ~\citep{bertsekas2012dynamic} but is expensive to compute exactly, we can devise a surrogate loss that approximates it. 
The operator $\Pi_{\Xi} \circ T_{\pi}$ is equivalent to the scaled projected gradient method $z_{t+1} = \Pi(z_t - \Xi^{-1} F(z_t))$, where $z_t = v_{\theta_t}$, and is minimum of the surrogate loss
\begin{align}
    \ell_t(\theta) &= \frac{1}{2}\|v_{\theta} -\left( v_{\theta_{t}} - \Xi^{-1} F(v_{\theta_t})\right)\|_{\Xi}^2\\ 
    &= \frac{1}{2}\|v_{\theta} - T_{\pi}(v_{\theta_t})\|_{\Xi}^2.
\end{align}
Note that if $\Xi = \id$ than this is the same surrogate loss \eqref{eq:surr} with $\eta=1$.

\subsection{The Linear Case}\label{sec:appendix-linear}

In this section we assume that predicted values are linear, $g(\theta) = \Phi \theta$, with some feature matrix
\[
    \Phi = \begin{bmatrix}
             &\phi^\top(s^1) \\
             &\vdots \\
             & \phi^\top(s^n)
            \end{bmatrix} \in \reals^{n \times d}.
\]
In this case we have that the operator $F$ is also linear in $\theta$
\[
    F(z) = F(\Phi \theta) = \Xi(\Phi \theta - \gamma P_\pi \Phi \theta  - r_\pi) = \Xi(\Phi - \gamma P_\pi \Phi) \theta  - \Xi r_\pi. 
\]
Since $g$ linear, there exists a preconditioning scheme that can generate a sequence of parameters $\{\theta_t\}_{t\in\nnumbers}$ such that $\{z_{t+1} = \Phi \theta_{t+1} =\Pi_{\Xi}(T_\pi(z_t)) \}$ and therefore is guaranteed to converge linearly.
This preconditioning scheme is simply the minimizer of the weighted least-squares surrogate loss $\ell_t(\theta)=\nicefrac{1}{2}\|\Phi \theta -T_{\pi}(\theta_t)\|^2_{\Xi}$,
\[
    \theta_{t+1} = \theta_t - (\Phi^\top \Xi \Phi)^{-1} \Phi^\top F(z_t). 
\]
If there is no unique minimizer then the pseudo-inverse $(\Phi^\top \Xi \Phi)^{\dag}$ can be used.
To approximate the iteration \citet{bertsekas2009projected} suggests using estimators that depend only on the history states $(s_1, s_2, \cdots,s_t, s_{t+1})$ and rewards $(r_1, \cdots, r_t)$ to approximate the conditioning matrix $\Phi^\top \Xi \Phi$ and gradient $\Phi^\top F(z_t)$.
The following estimators are proposed:
\begin{itemize}
    \item $\hat{D}_t = \frac{1}{t}\sum_{i=1}^t \phi(s_i)\phi^\top(s_i)$
    \item $\hat{C}_t = \frac{1}{t}\sum_{i=1}^t \phi(s_i)(\phi(s_i) - \gamma \phi(s_{i+1}))^\top$
    \item $\hat{r}_t=\frac{1}{t}\sum_{i=1}^t \phi(s_i)r_i$.
\end{itemize}
If the chain is fully mixed then the sampled states are drawn according to $\xi$ making the above estimators unbiased: 
\begin{align*}
    &\Phi^\top \Xi \Phi = \sum_i^n \xi^i \phi(s^i)\phi(s^i)^\top = \expect{\hat{D}_t},\\
    &\Phi^\top F(z) =  \sum_{i} \xi^i \phi(s^i)(\phi(s^i)^\top \theta  - \sum_{s'}(P_\pi)_{s^i,s} \phi(s)^\top \theta -r_{\pi}(s^i)) = \expect{\hat{C}_t \theta - \hat{r}_t}.
\end{align*}
Where $(P_\pi)_{s,s'}$ is the transition probability to state $s'$ when in state $s$.
Therefore the iteration
\[
\theta_{t+1} = \theta_t - (\hat{D}_t)^{-1}(\hat{C}_t\theta_t -\hat{r}_t), 
\]
eventually converges to exact deterministic iteration.
Interestingly, the stochastic version can be shown to be the minimmum of the surrogate
 \[
 \tilde{\ell}_t(\theta) = \frac{1}{2}\|\Phi \theta - \hat{T}_{\pi}(z_t)\|^2_{\hat{\Xi}}.
 \] 
 Where $\hat{T}_{\pi}(z) = \bar{r} + \gamma \bar{P}_\pi(z)$, with $\Xi$, $r$, and $P_{\pi}$ being estimated with the empirical distribution.
 More precisely, denoting the number of times that a state $s$ is visited in a trajectory of length $t$ as $n(s)$, and the number of times $s'$ is visited after $s$ as $n(s,s')$, then the empirical distribution over states is  $\hat{\xi}= (\nicefrac{n(s^1)}{t},\cdots, n(s^n)/t)$ and $\hat{\Xi}$ is the diagonal matrix with diagonal $\hat{\xi}$. 
 Similarly, the transition matrix is estimated by $(\bar{P}_{\pi})_{s,s'} = n(s,s')/n(s)$ if $n(s) > 0$, and arbitrary otherwise.
 For convenience, we write $\hat{\xi}(s^i) = \hat{\xi}^i$. 


\begin{proposition}\label{thm:bert-surr}
    The gradient of the surrogate loss $\tilde{\ell}_t(\theta)= \frac{1}{2}\|\Phi \theta - \hat{T}_{\pi}(z_t)\|^2_{\hat{\Xi}}$ is
    \[
        \nabla \tilde{\ell}_t(\theta) = \hat{C}_t\theta_t -\hat{r} + \hat{D}_t(\theta - \theta_t).
    \]
    If $\hat{D}_t$ is invertible then its minimum is
    \[
        \theta_{t+1} = \theta_t - (\hat{D}_t)^{-1}(\hat{C}_t \theta_t -\hat{r}_t).
    \]
\end{proposition}
\begin{proof}
    By definition of $\hat{\Xi}$, we have that $\hat{D}_t = \frac{1}{t}\sum_i \phi(s_i)\phi(s_i)^\top = \Phi^\top \hat{\Xi} \Phi$.
    Similarly, we have that $\hat{C}_t = \frac{1}{t}\sum_{i=1}^t \phi(s_i)(\phi(s_i) - \gamma \phi(s_{i+1}))^\top = \Phi^\top \hat{\Xi} \Phi - \gamma\frac{1}{t}\sum_{i=1}^t\phi(s_i)\phi(s_{i+1})^\top$.
    Where
    \begin{align}
        \frac{1}{t}\sum_{i=1}^t\phi(s_i)\phi(s_{i+1})^\top &= \frac{1}{t} \sum_{s,s'}n(s,s')\phi(s)\phi(s')^\top\\
        &= \frac{1}{t} \sum_{s}\phi(s)\sum_{s'} n(s,s') \phi(s')^\top\\
        &= \frac{1}{t} \sum_{s}n(s)\phi(s)\sum_{s'} \frac{n(s,s')}{n(s)} \phi(s')^\top\\
        &=  \frac{1}{t} \sum_{s}n(s)\phi(s)\sum_{s'} (\bar{P}_{\pi})_{s,s'}\phi(s')^\top\\
        &= \sum_s \hat{\xi}(s)\phi(s)\sum_{s'} (\bar{P}_{\pi})_{s,s'}\phi(s')^\top\\
        &= \Phi^\top \hat{\Xi}\bar{P}_\pi \Phi. 
    \end{align}
    Therefore $\hat{C}_t = \Phi^\top \hat{\Xi}(\Phi -\gamma \bar{P}_\pi \Phi)$.
    Similarly, $\hat{r}_t = \frac{1}{t}\sum_{i=1}^t \phi(s_i)r_i = \Phi^\top\Xi \bar{r}_\pi$.

    Now consider the weighted least squares loss
    \[
    \tilde{\ell}_t(\theta) = \frac{1}{2}\|\Phi \theta - \hat{T}_\pi(\Phi \theta_t)\|^2_{\hat{\Xi}}.
    \]
    The gradient of $\ell_t$ is
    \begin{align}
    \nabla \ell_t(\theta) &= \Phi^\top\hat{\Xi}(\Phi \theta - \hat{T}_{\pi}(\Phi \theta_t)) = \Phi^\top\hat{\Xi}(\Phi \theta - \gamma \bar{P}_{\pi}\Phi \theta_t  - \bar{r}_\pi)\\
    &=\Phi^\top\hat{\Xi}(\Phi \theta -\Phi \theta_t +\Phi \theta_t - \gamma \bar{P}_{\pi}\Phi \theta_t  - \bar{r}_\pi) \\ 
    &= \Phi^\top\hat{\Xi}(\Phi \theta -\Phi \theta_t) + \Phi^\top\hat{\Xi}(\Phi \theta_t - \gamma \bar{P}_{\pi}\Phi \theta_t  - \bar{r}_\pi)\\
    &= \hat{D}_t(\theta -\theta_t) + \hat{C}_t\theta_t -\hat{r}_t
   \end{align}
Setting the gradient to zero gives the minimum of $\tilde{\ell}_t$. 
\end{proof}

\subsection{The Nonlinear Case}
\begin{figure}
    \centering
    \includegraphics[width=\linewidth, scale=0.6]{figures/nonlinear_offline_exp_bg.pdf}
    \caption{\small Comparison of Bellman gap during training of TD(0) and surrogate methods in the offline full-batch setting with nonlinear function approximation in Ant (top) and HalfCheetah (bottom) environments, measured by gradient computations (left) and wallclock time (right). 
    The average Bellman gap \eqref{eq:approx-bellman-gap} across 20 runs along with 95\% confidence intervals are computed from a fixed training set. The red dashed line represents the lowest Bellman gap achieved by any of the algorithms. 
    }
    \label{fig:nonlinear_fullbatch_BG}
\end{figure}

\begin{figure}
    \centering
    \includegraphics[width=\linewidth, scale=0.6]{figures/nonlinear_offline_exp_ve.pdf}
    \caption{\small Comparison of average performance of TD(0) and surrogate methods in minimizing the value prediction error in the offline full-batch setting with nonlinear function approximation in Ant (top) and HalfCheetah (bottom) environments, measured by gradient computations (left) and wallclock time (right).
    The average value prediction error across 20 runs along with 95\% confidence intervals are computed from a fixed test set.
    The red dashed line represents the lowest value prediction error achieved by any of the algorithms.
    }
\label{fig:nonlinear_fullbatch_VE}
\end{figure}

In this section, we detail the surrogate algorithms used for the nonlinear setting, as outlined in Algorithms \ref{alg:surrogate_inner}, \ref{alg:double_sampling} and \ref{alg:thresholded_surr}.
As a complement to Section \ref{sec:experiments_PBE}, we compare the convergence properties of TD(0) with the surrogate methods under a an offline, full-batch setting, where the training data is pre-collected from the reference policy.
This experiment provides insight on two important fronts: how well can the methods solve the hidden monotone VI problem, and how the performance of the different algorithms are affected when there is no environment interaction.   

In our deep learning/nonlinear setting, the state space is too large to verify if we have solved the VI \eqref{eq:pbe-vi} by the $BE_{\text{gap}}$\eqref{eq:bellman-gap} or other metrics.
To test the effectiveness of our methods at solving this VI and minimizing PBE we consider an empirical approximation given a fixed batch of data (i.e. offline).
In this setting, $n$ tuples of state, reward and next state are collected $\{(s_t, r_t, s_{t+1})\}_{t=1}^n$ from a fixed reference policy $\pi$.
Using these observations we construct an empirical estimate of $BE$ \eqref{eq:pbe-loss-fixed-point},
\begin{equation}
    \widehat{BE}(\theta, \theta') = \frac{1}{n}\sum_{t=1}^n(v_{\theta}(s_t) - (r_t - \gamma v_{\theta'}(s_{t+1})))^2.
\end{equation}
We can then define the following fixed point problem:
find $\theta_{\ast}$ such that
\begin{equation}\label{eq:approx-pbe-loss-fixed-point}
    \theta_{\ast} \in \underset{\theta}{\arg\min} \, \widehat{BE}(\theta, \theta_\ast).
\end{equation}
Given this new empirical approximation of PBE we can define a gap analogous to \eqref{eq:bellman-gap}, 
\begin{equation}\label{eq:approx-bellman-gap}
    \widehat{BE}_{\text{gap}}(\theta) = \widehat{BE}(\theta, \theta) - \inf_{\theta'}\widehat{BE}(\theta', \theta).
\end{equation}
In practice we approximate the infimum in $\widehat{BE}$ by 500 steps of AdamW, with a learning rate of $1\times10^{-4}$ and exponential annealing at a decay rate of $\gamma = 0.995$.
In the offline setting we consider deterministic algorithms like in the min-max experiments of Section \ref{sec:experiments_minmax}, and therefore we do not  consider the double sampling algorithm (Algorithm \ref{alg:double_sampling}) but only full-batch versions of Algorithm \ref{alg:surrogate_inner} and Algorithm \ref{alg:thresholded_surr}.
We compare different algorithms with $\widehat{BE}_{\text{gap}}$ \eqref{eq:approx-bellman-gap}, giving a measure of how much the fixed point condition \eqref{eq:approx-pbe-loss-fixed-point} is violated. 
Given enough data this fixed point is a proxy for the fixed point \eqref{eq:pbe-loss-fixed-point}.

Since no environment interaction is required during training, we compare the performance of the methods in terms of gradient computations and wall-clock time, with the former being the dominant factor in computational cost, unlike the setting in Section \ref{sec:experiments_PBE} where environment interaction plays a larger role.

As shown in Figure~\ref{fig:nonlinear_fullbatch_BG}, increasing the number of inner loop steps significantly reduces the Bellman gap, suggesting that the model approaches the fixed point more effectively. Notably, the surrogate method with $10$ inner steps is more efficient than TD(0) in terms of gradient computations, while both inner steps = $10$ and $50$ outperform TD(0) in wall-clock time.

As presented in Figure~\ref{fig:nonlinear_fullbatch_VE}, we also evaluate the value prediction error on a test set to further assess the performance of the algorithms in the full batch and deterministic setting. In all cases, at least one variant of the surrogate method converges faster than TD(0), both in terms of gradient computations and wall-clock time. This demonstrates that, with an appropriate choice of inner loop steps and alpha threshold, the surrogate method not only outperforms TD(0) in terms of stability but also in computational efficiency.


\begin{figure}
    \centering
        \vspace{-5mm}
    \includegraphics[width=\linewidth, scale=0.6]{figures/experiments_LNN_new.pdf}
    \vspace{-7mm}
    \caption{
    \small
    Comparison of average performance of surrogate methods in minimizing the value prediction error for RL tasks with a 16-layer MLP in Ant as measured by outer loop iterations (left) and wallclock time (right).
    The average value prediction error across 20 runs along with 95\% confidence intervals are computed from a fixed test set.
    The red dashed line represents the lowest value prediction error achieved by any of the algorithms.
    }
    \label{fig:large_nonlinear}
        \vspace{-7mm}
\end{figure}
\subsubsection{Larger network}
In Figure \ref{fig:large_nonlinear} we report the average performance of surrogate methods using a 16 layer MLP. We observe similar runtimes and performance to the 2 layer MLP used in Section \ref{sec:non-linear}.

\begin{algorithm}[h]
\caption{Surrogate Method with Inner Loop}\label{alg:surrogate_inner}
\SetAlgoLined
\SetAlgoNoEnd
\KwIn{Reference policy $\pi$, initial value function parameters $\theta_0$, length of trajectories $N$, discount factor $\gamma$, learning rate $\eta$, number of inner steps $M$, number of outer loop iterations T, stop gradient function $\texttt{sg}(\cdot)$}

\For{$t = 0 \,\KwTo\, T$}{
    Collect a batch of trajectories $\{(s_i, a_i, r_i, s'_i)\}_{i=1}^N$ from $\pi$\\
    \For{each $(s_i, r_i, s'_i)$ in the batch}{
        Compute TD target: $y_i = r_i + \gamma \texttt{sg}(V_{\theta_t}(s'_i))$
    }
    Initialize inner loop $\theta$: $\theta_t^{(0)} = \theta_t$\\
    \For{$m = 0 \,\KwTo\, M-1$}{
        \For{each $(s_i, r_i, s'_i)$ in the batch}{
            Compute TD error: $\delta_i^{(m)} = V_{\theta_t^{(m)}}(s_i) - y_i$
        }
        Compute mean squared error: $\frac{1}{N}\sum_{i=1}^{N}(\delta_i^{(m)})^2$\\
        Update inner loop $\theta$: $\theta_t^{(m+1)} \leftarrow \theta_{t}^{(m)} - \eta \frac{1}{N} \sum_{i=1}^{N} \delta_i^{(m)} \nabla_{\theta} V_{\theta_t^{(m)}}(s_i)$
    }
    Update outer loop $\theta$: $\theta_{t+1} = \theta_{t}^{(M)}$
}
\end{algorithm}

\begin{algorithm}
\caption{Surrogate Method with Double Sampling}\label{alg:double_sampling}
\SetAlgoLined
\SetAlgoNoEnd
\KwIn{Reference policy $\pi$, initial value function parameters $\theta_0$, length of trajectories $N$, buffer size $K$, discount factor $\gamma$, learning rate $\eta$, number of inner steps $M$, number of outer loop iterations T, stop gradient function $\texttt{sg}(\cdot)$}
Collect $K$ batches of trajectories $\{(s_i, a_i, r_i, s'_i)\}_{i=1}^N$ from $\pi$ to construct a buffer\\
\For{$t = 0 \,\KwTo\, T$}{
    Collect a batch of trajectories $\{(s_i, a_i, r_i, s'_i)\}_{i=1}^N$ from $\pi$\\
    \For{each $(s_i, r_i, s'_i)$ in the batch}{
        Compute TD target: $y_i = r_i + \gamma \texttt{sg}(V_{\theta_t}(s'_i))$\\
        Compute TD error: $F(V_{\theta_t}(s_i)) = \texttt{sg}(V_{\theta_t}(s_i)) - y_i$
    }
    Initialize inner loop $\theta$: $\theta_t^{(0)} = \theta_t$\\
    \For{$m = 0 \,\KwTo\, M-1$}{
        Sample $m \mod K$-th batch $\{(s_j,a_j,r_j,s'_j)\}_{j=1}^N$ from the buffer\\
        \For{each $(s_i, r_i, s'_i)$ in the new batch}{
            Compute linearization term $l_1(\theta_t^{(m)}) = \langle F(V_{\theta_t}(s_i)), V_{\theta_{t}^{(m)}}(s_i) - \texttt{sg}(V_{\theta_{t}}(s_i)) \rangle$
        }
        \For{each $(s_j, r_j, s'_j)$ in the buffer batch}{
            Compute regularization term $l_2(\theta_t^{(m)}) = \frac{1}{2}\lVert V_{\theta_{t}^{(m)}}(s_j) - \texttt{sg}(V_{\theta_{t}}(s_j))\rVert_2^2$
        }
        Construct the surrogate loss function $L(\theta_t^{(m)}) = \frac{1}{N}(l_1(\theta_t^{(m)}) + l_2(\theta_t^{(m)}))$\\
        Update inner loop $\theta_{t}^{(m)}$ by minimizing the surrogate loss: $\theta_{t}^{(m+1)} \leftarrow \theta_{t}^{(m)} - \eta \nabla_{\theta}L(\theta_t^{(m)})$
    }
    Update outer loop $\theta$: $\theta_{t+1} = \theta_{t}^{(M)}$
}
\end{algorithm}

\begin{algorithm}[H]
\caption{Thresholded Surrogate Method}\label{alg:thresholded_surr}
\SetAlgoLined
\SetAlgoNoEnd
\KwIn{Reference policy $\pi$, initial value function parameters $\theta_0$, length of trajectories $N$, discount factor $\gamma$, learning rate $\eta$, number of inner steps $M$, number of outer loop iterations T, stop gradient function $\texttt{sg}(\cdot)$, $\alpha^2\in(0,1)$}
\For{$t = 0 \,\KwTo\, T$}{
    Collect a batch of trajectories $\{(s_i, a_i, r_i, s'_i)\}_{i=1}^N$ from $\pi$\\
    \For{each $(s_i, r_i, s'_i)$ in the batch}{
        Compute TD target: $y_i = r_i + \gamma \texttt{sg}(V_{\theta_t}(s'_i))$
    }
    Initialize inner loop $\theta$: $\theta_t^{(0)} = \theta_t$\\
    Compute initial surrogate loss: $l_t(\theta_t)=\frac{1}{N}\sum_{i=1}^{N}\lVert V_{\theta_t}(s_i)-y_i\rVert_2^2$\\
    \While{$l_t(\theta_t^{(m)})/l_t(\theta_t)\geq\alpha^2$ and $m\leq M$}{
        \For{each $(s_i, r_i, s'_i)$ in the batch}{
            Compute TD error: $\delta_i^{(m)} = V_{\theta_t^{(m)}}(s_i) - y_i$
        }
        Compute mean squared error: $\frac{1}{N}\sum_{i=1}^{N}(\delta_i^{(m)})^2$\\
        Update inner loop $\theta$: $\theta_t^{(m+1)} \leftarrow \theta_{t}^{(m)} - \eta \frac{1}{N} \sum_{i=1}^{N} \delta_i^{(m)} \nabla_{\theta} V_{\theta_t^{(m)}}(s_i)$\\
        Update $m: m\leftarrow m+1$
    }
    Update outer loop $\theta$: $\theta_{t+1} = \theta_{t}^{(M)}$
}
\end{algorithm}

\section{Applications to Performative Prediction}

Minimizing Projected Bellman error can be viewed as a special case of finding a stable point in performative prediction~\citep{perdomo2020performative}.
In performative prediction or more generally optimization under decision dependent distributions~\citep{drusvyatskiy2023stochastic}, the model predictions $z=g(\theta)$ influence the targets $y$.
With moving targets it is of interest to find model predictions that are performatively stable; such models are optimal under the shifting targets and are stable under repeated applications of empirical risk minimization.

For simplicity, we stay within the finite dimensional case where a model's prediction function is given by the $n$-dimensional vector, $z = g(\theta) \in \reals^n$.
Where each prediction is $z^i = g^i(\theta) = h(x_i, \theta)$, for some feature vector $x_i$ and fixed model architecture $h$.
Given an error function $\ell(z,y)$, quantifying error between predictions $z \in \reals^n$ and targets $y \in \reals^n$, the performatively stable models satisfy
\begin{align}\label{eq:pp-stable}
    \theta_\ast \in \underset{\theta}{\arg \min} \, \expect[y\sim D(g(\theta_\ast))]{\ell(g(\theta),y)}. 
\end{align}
The prediction dependent shift in targets is represented by the distribution $D(\cdot)$ being a function of the model predictions $g(\theta)$.
Note that we have used the formulation of performative prediction with respect to model predictions~\citep{mofakhami2023performative} instead of parameters.

The problem of minimizng projected Bellman error is a special case of \eqref{eq:pp-stable} with $\ell(z,y) = ||z-y||^2_{\Xi}$ and $D(z)$ is the distribution with full weight on $T_{\pi}(z)$. 

The problem of finding a performatively stable point can more generally be expressed as the following decision-dependent optimization problem: 
\begin{align}
    \theta_\ast \in \underset{\theta }{\arg \min } f_{g(\theta_\ast)}(g(\theta)).
\end{align}
Or equivalently a constrained optimization problem with constraint $\mathcal{Z} = \closure\{g(\theta) : \theta \in \reals^d\}$:
\begin{align*}
    \theta_\ast &\in \underset{\theta }{\arg \min } f_{g(\theta_\ast)}(g(\theta)) \Leftrightarrow z_\ast \in \underset{ z \in \mathcal{Z}}{\arg \min } f_{z_\ast}(z).
\end{align*}
If $Z$ is closed convex and $f_{z'}(z)$ is convex and differentiable with respect to $z$, then the above constrainted formulation is equivalent to the fixed point problem:
\begin{align}\label{eq:pp-fixed-point}
    z_{\ast} = \Pi \circ T (z_\ast) = \Pi(z_\ast -\eta \nabla f_{z_\ast}(z_\ast)),
\end{align}
where we denote the gradient of $f_{z'}(z)$ with respect to $z$ as simply $\nabla f_{z'}(z)$.
As mentioned in Section \ref{sec:appendix-PBE}, \citet{bertsekas2009projected} showed that solving \eqref{eq:pp-fixed-point} is equivalent to solving a strongly monotone and smooth VI if either $T$ or $\Pi \circ T$ is a contraction.
Therefore, under these contraction conditions \eqref{eq:pp-fixed-point} is a hidden strongly monotone and smooth VI problem of the form \eqref{eq:vi}! 
Thankfully, \citep{drusvyatskiy2023stochastic, cheng2020online} showed that $\Pi \circ T$ is a contraction if $f_{z'}(z)$ is $\mu$ strongly convex with respect to $z$ and $\nabla f_{z'}(z)$ is $\beta$-Lipschitz with respect to $z'$ and if $\nicefrac{\beta}{\mu} < 1$.
Therefore, under standard performative prediction assumptions such as those found in \citet{drusvyatskiy2023stochastic}, the performatively stable point \eqref{eq:pp-stable} is equivalent to solving a hidden strongly monotone and smooth VI problem where our surrogate loss approach (Algorithm \eqref{alg:surr}) applies and converges to the stable point under Theorems (\ref{thm:det},\ref{thm:stochastic},\ref{thm:contraction-error-limit}). 

\end{document}
